% count/count.tex
% mainfile: ../perfbook.tex
% SPDX-License-Identifier: CC-BY-SA-3.0

\QuickQuizChapter{chp:Counting}{计数}{qqzcount}
%
\Epigraph{简单如1、2、3！}{佚名}

计数也许是计算机能做的事情里，最简单也最自然的工作了。
不过在大型共享内存多处理器系统上，想要高效且可扩展地计数，
仍然颇具挑战性。不仅如此，计数这一概念的简洁性让我们可以
在探索并发的基本问题时，无须被繁复的数据结构或精巧的同步原语干扰。
因此，计数是并行编程的绝佳切入对象。

本章涵盖了许多适用于特殊情况的简单、快速且可扩展的计数算法。
但首先，让我们来看看你对并发计数了解多少。

\EQuickQuiz{
	为什么高效且可扩展的计数会很困难？
	毕竟，计算机拥有专门用于计数的硬件！
}\EQuickQuizAnswer{
	因为直接的计数算法——例如对共享counter进行atomic操作——要么
	速度慢且扩展性差，要么不准确，这将在
	\cref{sec:count:Why Isn't Concurrent Counting Trivial?}中说明。
}\EQuickQuizEnd

\EQuickQuiz{
	{\bfseries 网络数据包计数问题。}
	假设你需要收集已发送和接收的网络数据包数量的统计信息。
	数据包可能由系统上的任何CPU发送或接收。
	进一步假设你的系统每个CPU每秒能处理数百万个数据包，
	而系统监控程序每五秒读取一次计数值。
	你将如何实现这个counter？
}\EQuickQuizAnswer{
	提示：
	更新counter的操作必须极快，但由于counter大约每五百万次更新
	才被读取一次，读取counter的操作可以相当慢。
	此外，读取的值通常不需要非常精确——毕竟，既然counter
	每毫秒更新一千次，我们应该能够接受与``真实值''相差
	几千次的值，不管在这个语境下``真实值''意味着什么。
	但是，读取的值应随时间保持大致相同的绝对误差。
	例如，当计数值在百万级别时，1\pct\ 的误差可能完全可以接受，
	但当计数值达到万亿级别时，这可能是完全不可接受的。
	参见\cref{sec:count:Statistical Counters}。
}\EQuickQuizEnd

\QuickQuizLabel{\QcountQstatcnt}

\EQuickQuiz{
	{\bfseries 近似结构分配限制问题。}
	假设你需要维护已分配结构数量的计数，以便在使用中的结构数量
	超过限制（比如10,000）时拒绝任何分配。
	进一步假设这些结构是短生命周期的，限制很少被超过，
	且``宽松的''近似限制是可以接受的。
}\EQuickQuizAnswer{
	提示：
	更新counter的操作必须再次极快，但每次增加counter时
	都要读取它。
	然而，读取的值不需要精确，
	\emph{除了}它必须能大致区分低于限制的值和大于或
	等于限制的值。
	参见\cref{sec:count:Approximate Limit Counters}。
}\EQuickQuizEnd

\QuickQuizLabel{\QcountQapproxcnt}

\EQuickQuiz{
	{\bfseries 精确结构分配限制问题。}
	假设你需要维护已分配结构数量的计数，以便在使用中的结构数量
	超过精确限制（同样假设为10,000）时拒绝任何分配。
	进一步假设这些结构是短生命周期的，
	限制很少被超过，几乎总是至少有一个结构在使用中，
	并且还需要精确知道counter何时到达零——例如，
	为了释放某些只在至少有一个结构在使用时才需要的内存。
}\EQuickQuizAnswer{
	提示：
	更新counter的操作必须再次极快，但每次增加counter时
	都要读取它。
	然而，读取的值不需要精确，
	\emph{除了}它绝对必须完美地区分处于限制值与零之间的值
	和小于或等于零或大于或等于限制值的值。
	参见\cref{sec:count:Exact Limit Counters}。
}\EQuickQuizEnd

\QuickQuizLabel{\QcountQexactcnt}

\EQuickQuiz{
	{\bfseries 可移除I/O设备访问计数问题。}
	假设你需要在一个高频使用的可移除大容量存储设备上
	维护一个\IX{reference count}，
	以便能够告诉用户何时可以安全移除该设备。
	按照惯例，用户表示希望移除设备，
	系统告知用户何时可以安全执行此操作。
}\EQuickQuizAnswer{
	提示：
	再次，更新counter的操作必须极快且可扩展，
	以避免减慢I/O操作，
	但由于counter仅在用户希望移除设备时才被读取，
	counter的读取操作可以极慢。
	此外，除非用户已经表示希望移除设备，
	否则根本不需要读取counter。
	另外，读取的值不需要精确，
	\emph{除了}它绝对必须完美地区分非零值和零值，
	而且仅在设备处于移除过程中时才需如此。
	但是，一旦读取到零值，就必须采取行动
	将该值保持为零，直到已采取某些措施
	阻止后续线程获取对正在移除的设备的访问。
	参见\cref{sec:count:Applying Exact Limit Counters}。
}\EQuickQuizEnd

\QuickQuizLabel{\QcountQIOcnt}

\Cref{sec:count:Why Isn't Concurrent Counting Trivial?}
提出问题：为什么在多核系统上的计数问题不简单？
\Cref{sec:count:Statistical Counters,sec:count:Approximate Limit Counters}
分别探索了网络数据包计数问题和近似结构分配上限问题。
\Cref{sec:count:Exact Limit Counters}
研究精确结构分配上限问题。
最后，\cref{sec:count:Parallel Counting Discussion}
对本章进行总结，同时展示各种方法的性能测量结果。

\Cref{sec:count:Why Isn't Concurrent Counting Trivial?,%
sec:count:Statistical Counters}
都是介绍性的内容，之后的内容更为高级。

\section{为什么并发计数不可小看？}
\label{sec:count:Why Isn't Concurrent Counting Trivial?}
%
\epigraph{追求简洁，但不要轻信它。}{Alfred North Whitehead}

让我们从简单的东西开始，例如
\cref{lst:count:Just Count!}（\path{count_nonatomic.c}）中展示的
直接使用算术运算的方式。
\begin{fcvref}[ln:count:count_nonatomic:inc-read]
这里，我们在\clnref{counter}处有一个counter，在\clnref{inc}处
递增它，并在\clnref{read}处读取其值。
还有什么比这更简单的呢？
\end{fcvref}

\QuickQuiz{
	比这更简单的一件事是使用\co{++}而不是那串
	\co{READ_ONCE()}和\co{WRITE_ONCE()}的拼接。
	为什么要多打那么多字？
}\QuickQuizAnswer{
	参见\cref{sec:toolsoftrade:Shared-Variable Shenanigans}
	（位于\cpageref{sec:toolsoftrade:Shared-Variable Shenanigans}），
	了解更多关于编译器如何制造麻烦的信息，
	以及\co{READ_ONCE()}和\co{WRITE_ONCE()}如何避免
	这些麻烦。
}\QuickQuizEnd

\begin{listing}
\input{CodeSamples/count/count_nonatomic@inc-read.fcv}
\caption{直接计数！}
\label{lst:count:Just Count!}
\end{listing}

当计数器不停读取却又几乎不增加计数时，这种方法快得令人惊叹。
在小型系统上，性能也非常出色。

不过让人扫兴的是，这种方法可能丢失计数。
在我的六核x86笔记本电脑上，一次短暂运行调用了\co{inc_count()}
285,824,000次，但counter的最终值仅为35,385,525。
虽然近似值在计算领域确实占有一定地位，但丢失87\pct\ 的计数
未免过于夸张了。

\QuickQuizSeries{%
\QuickQuizB{
	难道一个聪明的编译器不能证明
	\cref{lst:count:Just Count!}的
	\clnrefr{ln:count:count_nonatomic:inc-read:inc}
	等价于\co{++}运算符，从而生成一条x86加法到内存的指令吗？
	而且CPU缓存不会使其成为atomic的吗？
}\QuickQuizAnswerB{
	虽然\co{++}运算符\emph{可能}是atomic的，但除非将其应用于
	C11的\co{_Atomic}变量，否则没有要求它必须如此。
	事实上，在没有\co{_Atomic}的情况下，\GCC\ 通常
	选择将值加载到寄存器中，递增寄存器，然后将值存储回内存，
	这显然不是atomic的。

	此外，请注意\co{READ_ONCE()}和\co{WRITE_ONCE()}中的
	volatile强制类型转换，它们告诉编译器该位置很可能是一个
	MMIO设备寄存器。
	因为MMIO寄存器不被缓存，编译器假设递增操作是atomic的
	将是不明智的。
}\QuickQuizEndB
%
\QuickQuizE{
	失败次数的8位精度表明你确实测试了这个程序。
	为什么需要测试如此简单的程序，
	特别是当这个bug通过检查代码就能轻易发现时？
}\QuickQuizAnswerE{
	不仅几乎没有简单的并行程序，而且大多数时候我
	也不敢确定有多少简单的串行程序。

	无论程序多小多简单，如果你没有测试过，
	它就不能工作。
	即使你测试过了，墨菲定律也告诉我们至少还有
	一些bug潜伏着。

	此外，虽然正确性证明确实有其用武之地，但它们永远
	无法替代测试，包括此处使用的\path{counttorture.h}
	测试框架。
	毕竟，证明的可靠性取决于其所基于的假设。
	最后，证明本身可能和程序一样存在bug！
}\QuickQuizEndE
}

\begin{listing}
\input{CodeSamples/count/count_atomic@inc-read.fcv}
\caption{原子计数！}
\label{lst:count:Just Count Atomically!}
\end{listing}

\begin{figure}
\centering
\resizebox{2.5in}{!}{\includegraphics{CodeSamples/count/atomic_hps.pdf}}
\caption{x86上的Atomic递增可扩展性}
\label{fig:count:Atomic Increment Scalability on x86}
\end{figure}

精确计数的最直接的办法是使用\IX{atomic}操作，
如\cref{lst:count:Just Count Atomically!}（\path{count_atomic.c}）
所示。
\begin{fcvref}[ln:count:count_atomic:inc-read]
\Clnref{counter}定义了一个atomic变量，
\clnref{inc}对其进行atomic递增，
\clnref{read}读取其值。
\end{fcvref}
因为都是atomic操作，所以计数非常精确。
不过，程序执行得要慢一些：
在我的六核x86笔记本电脑上，即使只有单个线程（thread）在递增，
它也比非atomic递增慢了二十多倍。\footnote{
	有趣的是，非atomic地递增counter会比atomic地递增
	counter更快地推进counter值。
	当然，如果你的唯一目标是让counter快速增长，
	更简单的方法是直接给counter赋一个大值。
	尽管如此，使用适当放松的正确性概念来获得
	更好的性能和可扩展性的算法很可能有其用武之地
	~\cite{Andrews91textbook,Arcangeli03,10.5555/3241639.3241645,DavidUngar2011unsync}。}

这种糟糕的性能并不让人意外，鉴于
\cref{chp:Hardware and its Habits}中的讨论，
atomic递增的性能随CPU和线程数量的增加而下降也不足为奇，
如\cref{fig:count:Atomic Increment Scalability on x86}所示。
在这张图中，x轴上的水平虚线代表
完美可扩展算法所能达到的理想性能：
在这样的算法中，CPU和线程数目的增长带来的开销
与单线程程序中的开销相同。
对单个全局变量的atomic递增显然远非理想，
随着CPU个数的增加，性能差距达到数个数量级。

\QuickQuizSeries{%
\QuickQuizB{
	为什么x轴上的水平虚线没有在$x=1$处与对角线相交？
}\QuickQuizAnswerB{
	因为atomic操作的开销。
	x轴上的虚线表示单次\emph{非atomic}递增的开销。
	毕竟，一个\emph{理想的}算法不仅应该线性扩展，
	还应该与单线程代码相比不产生任何性能损失。

	这种理想主义标准可能看起来很严格，但如果它对
	\ppl{Linus}{Torvalds}来说足够好，那对你来说也足够好。
}\QuickQuizEndB
%
\QuickQuizE{
	但atomic递增还是挺快的。
	而且在紧密循环中递增单个变量听起来
	相当不切实际，毕竟程序的大部分执行时间应该
	用于实际工作，而不是记录已完成的工作！
	为什么我要关心让它更快？
}\QuickQuizAnswerE{
	在许多情况下，atomic递增实际上足够快。
	在这些情况下，你完全应该使用atomic递增。

	话虽如此，在许多实际场景中确实需要
	更复杂的计数算法。
	这种场景的典型例子是在高度优化的网络栈中
	计数数据包和字节，在大型多处理器上很容易
	发现大量执行时间花在了这类记账任务上。
	事实上，一些研究者定义了``超越单线程的配置''
	（COST），发现令人惊讶的是，通常需要大量线程
	才能超越单线程的性能~\cite{FrankMcSherry2015COST}。
	换句话说，这不是一个微不足道的问题，
	如果不付出努力，它不太可能自行消失。

	简而言之，正如本章开头所指出的，计数
	为共享内存并行程序中遇到的问题提供了
	极好的视角。
}\QuickQuizEndE
}

\begin{figure}
\centering
\resizebox{3in}{!}{\includegraphics{count/GlobalInc}}
\caption{全局Atomic递增的数据流}
\label{fig:count:Data Flow For Global Atomic Increment}
\end{figure}

\begin{figure}
\centering
\resizebox{3.2in}{!}{\includegraphics{cartoons/r-2014-One-one-thousand}}
\caption{等待计数}
\ContributedBy{fig:count:Waiting to Count}{Melissa Broussard}
\end{figure}

\cref{fig:count:Data Flow For Global Atomic Increment}提供了
另一种看全局atomic递增的视角。
为了让每个CPU得到机会递增一个指定的全局变量，
包含该变量的\IX{cache line}需要在所有CPU之间传播，
如红色箭头所示。
由于光速有限和原子大小非零，这种传播相当耗时，
这些物理定律在\cref{sec:cpu:Hardware Free Lunch?}中已经讨论过。
这种缓慢的传播反过来导致了
\cref{fig:count:Atomic Increment Scalability on x86}中
所见的糟糕性能，
\cref{fig:count:Waiting to Count}形象地描绘了这种情况。
以下章节讨论高性能计数，这种方法可以避免此类传播所固有的延迟。

\QuickQuiz{
	但是为什么CPU设计者不能简单地将加法操作发送到
	数据所在位置，从而避免需要流转包含被递增的全局变量的
	cache line呢？
}\QuickQuizAnswer{
	在某些情况下可能确实能做到这一点。
	但是，存在一些复杂因素：
	\begin{enumerate}
	\item	如果需要变量的值，那么线程将被迫等待
		操作被发送到数据所在位置，然后等待结果
		被返回。
	\item	如果atomic递增必须相对于之前和/或之后的操作
		有序，那么线程将被迫等待操作被发送到
		数据所在位置，并等待操作完成的指示被返回。
	\item	在CPU之间传输操作可能需要系统互连中的
		更多线路，这将消耗更多的芯片面积和更多的电能。
	\end{enumerate}
	但如果前两个条件都不成立呢？
	那么你应该仔细考虑\cref{sec:count:Statistical Counters}中
	讨论的算法，它们在通用硬件上实现了接近理想的性能。

\begin{figure}
\centering
\resizebox{3in}{!}{\includegraphics{count/GlobalTreeInc}}
\caption{全局组合树Atomic递增的数据流}
\label{fig:count:Data Flow For Global Combining-Tree Atomic Increment}
\end{figure}

	如果前两个条件中的一个或两个成立，硬件改进
	还是有\emph{一些}希望的。
	可以设想硬件实现一个组合树，
	这样来自多个CPU的递增请求在到达硬件时
	被合并为单次加法。
	硬件还可以对请求施加顺序，从而
	向每个CPU返回与其特定atomic递增对应的返回值。
	这导致指令延迟为$\O{\log N}$，
	其中$N$是CPU数量，如
	\cref{fig:count:Data Flow For Global Combining-Tree Atomic Increment}
	所示。
	据说一些具有这种硬件优化的CPU在1990年代曾投入生产使用，
	并在2011年开始重新出现。\footnote{
		第一个发送可引用参考文献来确认或推翻此陈述的读者
		将在本脚注的未来版本中被提及。}

	相比\cref{fig:count:Data Flow For Global Atomic Increment}
	中所示的当前硬件的$\O{N}$性能，
	这是一个巨大的改进，
	如果三维制造等创新证明可行的话，
	硬件延迟可能进一步降低。
	尽管如此，我们将看到在一些重要的特殊情况下，
	软件能做得\emph{好得多}。
}\QuickQuizEnd

\section{统计计数器（statistical counter）}
\label{sec:count:Statistical Counters}
%
\epigraph{事实是固执的东西，但统计数字是可以揉捏的。}
	 {Mark Twain}

本节涵盖使用statistical counters的常见场景，其中计数值
被极其频繁地更新，而很少甚至从不被读取。
这些方法将用于解决\QuickQuizRef{\QcountQstatcnt}中提出的
网络数据包计数问题。

\subsection{设计}

Statistical counting一般以每个线程一个counter的方式处理
（或者在内核中运行时，每个CPU一个），这样每个线程只更新自己的counter，
正如\cref{sec:toolsoftrade:Per-CPU Variables}
（\cpageref{sec:toolsoftrade:Per-CPU Variables}）中所预示的那样。
根据加法的交换律和结合律，总的计数值就是所有线程counter值的简单相加。
这是Data Ownership模式的一个例子，将在
\cref{sec:SMPdesign:Data Ownership}
（\cpageref{sec:SMPdesign:Data Ownership}）中介绍。

\QuickQuiz{
	但是C语言``整数''大小有限这个事实不会使事情复杂化吗？
}\QuickQuizAnswer{
	不会，因为模加法仍然满足交换律和结合律。
	至少在你使用无符号整数的情况下是这样。
	回忆一下，在C标准中，有符号整数的溢出导致未定义行为，
	尽管事实上如今在溢出时做任何其他事情（而非回绕）
	的机器相当罕见。
	不幸的是，编译器经常执行假设有符号整数不会溢出的优化，
	因此如果你的代码允许有符号整数溢出，即使在现代的
	二进制补码硬件上也可能遇到麻烦。

	话虽如此，当试图从32位per-thread counter中
	收集（比方说）64位的和时，确实可能产生额外的复杂性。
	处理这种额外复杂性留作读者的练习，
	本章后面介绍的一些技术可能会很有帮助。
}\QuickQuizEnd

\subsection{基于数组的实现}
\label{sec:count:Array-Based Implementation}

实现per-thread变量的一种方法是分配一个数组，数组每个元素
对应一个线程（假设数组已按cache line对齐并且填充过了，以避免false sharing）。

\QuickQuiz{
	数组？
	但这不是限制了线程数量吗？
}\QuickQuizAnswer{
	确实可能，在这个玩具实现中确实如此。
	但要想出一个允许任意数量线程的替代实现
	并不困难，例如使用C11的\co{_Thread_local}功能，
	如\cref{sec:count:Per-Thread-Variable-Based Implementation}
	所示。
}\QuickQuizEnd

\begin{listing}
\input{CodeSamples/count/count_stat@inc-read.fcv}
\caption{基于数组的Per-Thread Statistical Counters}
\label{lst:count:Array-Based Per-Thread Statistical Counters}
\end{listing}

这样的数组可以被封装成per-thread原语，如
\cref{lst:count:Array-Based Per-Thread Statistical Counters}
（\path{count_stat.c}）所示。
\begin{fcvref}[ln:count:count_stat:inc-read]
\Clnref{define}定义了一个包含per-thread counter集合的数组，
类型为\co{unsigned long}，创造性地命名为\co{counter}。

\Clnrefrange{inc:b}{inc:e}
展示了一个递增counter的函数，使用
\apipf{__get_thread_var()}原语来定位当前运行线程
在\co{counter}数组中的元素。
因为该元素只能由对应的线程修改，
非atomic递增即可满足。
然而，此代码使用\co{WRITE_ONCE()}来防止破坏性的
编译器优化。
举个例子，编译器有权将即将被存储的位置
用作临时存储，从而在执行所需存储之前写入
实际上是垃圾的内容。
这当然可能会让试图读取计数值的代码感到困惑。
使用\co{WRITE_ONCE()}可以防止这种优化及其他类似优化。

\QuickQuiz{
	\GCC\ 还可能应用什么其他恶意优化？
}\QuickQuizAnswer{
	参见\cref{sec:toolsoftrade:Shared-Variable Shenanigans,%
	sec:memorder:Compile-Time Consternation}
	获取更多信息。
	一种恶意优化是对\co{read_count()}函数的连续调用
	应用公共子表达式消除，这可能会让期望从连续调用中
	获得不同返回值的代码感到意外。
}\QuickQuizEnd

\Clnrefrange{read:b}{read:e}
展示了一个读取counter聚合值的函数，
使用\apipf{for_each_thread()}原语遍历当前运行线程的列表，
并使用\apipf{per_thread()}原语获取指定线程的counter。
此代码还使用\co{READ_ONCE()}来确保编译器不会
将这些加载操作优化掉。
举一个例子，一对连续的\co{read_count()}调用可能被内联，
一个大胆的优化器可能注意到相同的位置被求和了两次，
从而错误地得出结论：只求和一次并使用结果两次
是个好主意。
这种优化可能会让期望后续\co{read_count()}调用
能反映其他线程活动的人感到沮丧。
使用\co{READ_ONCE()}可以防止这种优化及其他类似优化。
\end{fcvref}

\QuickQuizSeries{%
\QuickQuizB{
	\cref{lst:count:Array-Based Per-Thread Statistical Counters}
	中的per-thread \co{counter}变量是如何初始化的？
}\QuickQuizAnswerB{
	C标准规定全局变量的初始值为零（除非显式初始化），
	因此隐式地将\co{counter}的所有实例初始化为零。
	此外，在用户只对statistical counters连续读取之间的
	差值感兴趣的常见情况下，初始值是无关紧要的。
}\QuickQuizEndB
%
\QuickQuizE{
	\cref{lst:count:Array-Based Per-Thread Statistical Counters}
	中的代码如何支持多于一个counter？
}\QuickQuizAnswerE{
	确实，这个玩具示例不支持多于一个counter。
	修改它使其能够提供多个counter留作读者的练习。
}\QuickQuizEndE
}

\begin{figure}
\centering
\resizebox{3in}{!}{\includegraphics{count/PerThreadInc}}
\caption{Per-Thread递增的数据流}
\label{fig:count:Data Flow For Per-Thread Increment}
\end{figure}

该方法随着调用\co{inc_count()}的更新者线程增加而线性扩展。
如\cref{fig:count:Data Flow For Per-Thread Increment}中每个CPU上的
绿色箭头所示，原因是每个CPU可以快速递增
自己线程的变量，不再需要代价昂贵的跨系统通信。
换句话说，per-CPU counter的更新不会像单个全局counter的
atomic递增那样受到慢光和大原子的困扰。\footnote{
	同样，参见\cref{sec:cpu:Hardware Free Lunch?}。}
因此，本节解决了本章开头提出的网络数据包计数问题。

\QuickQuiz{
	读取操作需要时间来汇总per-thread的值，
	在此期间counter很可能正在变化。
	这意味着
	\cref{lst:count:Array-Based Per-Thread Statistical Counters}
	中\co{read_count()}返回的值不一定精确。
	假设counter以每单位时间$r$次的速率递增，
	而\co{read_count()}的执行消耗$\Delta$个时间单位。
	返回值的期望误差是多少？
}\QuickQuizAnswer{
	让我们先做最坏情况分析，然后做一个不那么保守的分析。

	在最坏情况下，读取操作立即完成，
	但在返回之前延迟了$\Delta$个时间单位，
	此时最坏情况误差简单地为$r \Delta$。

	这种最坏情况行为相当不可能，因此让我们考虑
	从$N$个counter中的每个读取均匀分布在时间段$\Delta$内的情况。
	在$N$次读取之间将有$N+1$个持续时间为$\frac{\Delta}{N+1}$
	的间隔。
	递增速率$r$预计均匀分布在$N$个counter上，
	每个counter每单位时间$\frac{r}{N}$次递增。
	从最后一个线程的counter读取后的延迟引起的误差
	为$\frac{r \Delta}{N \left( N + 1 \right)}$，
	倒数第二个线程的counter为
	$\frac{2 r \Delta}{N \left( N + 1 \right)}$，
	倒数第三个为
	$\frac{3 r \Delta}{N \left( N + 1 \right)}$，
	依此类推。
	总误差由从每个线程counter读取产生的误差之和给出：

	\begin{equation}
		\frac{r \Delta}{N \left( N + 1 \right)}
			\sum_{i = 1}^N i
	\end{equation}

	将求和以闭式形式表示：

	\begin{equation}
		\frac{r \Delta}{N \left( N + 1 \right)}
			\frac{N \left( N + 1 \right)}{2}
	\end{equation}

	约分得到直觉上预期的结果：

	\begin{equation}
		\frac{r \Delta}{2}
	\label{eq:count:CounterErrorAverage}
	\end{equation}

	重要的是要记住，当调用者执行利用读取操作
	返回的计数值的代码时，误差会继续累积。
	例如，如果调用者花费时间$t$执行基于返回计数值
	的某些计算，最坏情况误差将增加到$r \left(\Delta + t\right)$。

	期望误差将类似地增加到：

	\begin{equation}
		r \left( \frac{\Delta}{2} + t \right)
	\end{equation}

	当然，有时counter在读取操作期间继续递增
	是不可接受的。
	\Cref{sec:count:Applying Exact Limit Counters}
	讨论了处理这种情况的方法。

	到目前为止，我们一直在考虑一个只增不减的counter。
	如果counter值以每单位时间$r$次的速率变化，
	但可以是任一方向，我们应该期望误差会减少。
	然而，最坏情况不变，因为虽然counter\emph{可以}
	向任一方向移动，但最坏情况是读取操作立即完成，
	然后延迟$\Delta$个时间单位，在此期间
	counter值的所有变化都朝同一方向，
	再次给出$r \Delta$的绝对误差。

	有多种方法可以基于各种关于递增和递减模式的
	假设来计算平均误差。
	为简单起见，让我们假设$f$比例的操作是递减，
	而感兴趣的误差是偏离counter长期趋势线的偏差。
	在此假设下，如果$f$小于或等于0.5，
	每次递减将被一次递增抵消，因此
	$2f$的操作将相互抵消，留下
	$1-2f$的操作是未抵消的递增。
	另一方面，如果$f$大于0.5，$1-f$的
	递减被递增抵消，因此counter在负方向移动
	$-1+2\left(1-f\right)$，
	简化为$1-2f$，因此counter在两种情况下
	平均每次操作移动$1-2f$。
	因此，counter的长期移动由$\left( 1-2f \right) r$给出。
	代入\cref{eq:count:CounterErrorAverage}得到：

	\begin{equation}
		\frac{\left( 1 - 2 f \right) r \Delta}{2}
	\end{equation}

	除上述所有内容外，在statistical counters的大多数用途中，
	\co{read_count()}返回值中的误差是无关紧要的。
	这种无关紧要是因为\co{read_count()}执行所需的时间
	通常远小于\co{read_count()}连续调用之间的时间间隔。
}\QuickQuizEnd

\QuickQuizLabel{\StatisticalCounterAccuracy}

然而，许多实现提供了更廉价的per-thread数据机制，
且不受数组大小的限制。
下一节将展示这样的方法。

\subsection{基于per-thread变量的实现}
\label{sec:count:Per-Thread-Variable-Based Implementation}

C语言自C11起提供了\apic{_Thread_local}存储类，
用于提供per-thread存储。\footnote{
	\GCC\ 提供了自己的\apig{__thread}存储类，
	在本书以前的版本中使用。
	使用\GCC\ 时，这两种指定thread-local变量的方法
	是可以互换的。}
如\cref{lst:count:Per-Thread Statistical Counters}
（\path{count_end.c}）所示，
可以使用它来实现一个statistical counter，这个counter不仅能扩展，
而且相对于简单的非atomic递增来说，对递增方几乎不产生性能损失。

\begin{listing}
\input{CodeSamples/count/count_end@whole.fcv}
\caption{Per-Thread Statistical Counters}
\label{lst:count:Per-Thread Statistical Counters}
\end{listing}

\begin{fcvref}[ln:count:count_end:whole]
\Clnrefrange{var:b}{var:e}定义了所需的变量：
\co{counter}是per-thread counter变量，
\co{counterp[]}数组允许线程访问彼此的counter，
\co{finalcount}在各个线程退出时将计数值累加到总和里，
\co{final_mutex}用于协调累加计数总值的线程和正在退出的线程。
\end{fcvref}

\QuickQuiz{
	\cref{lst:count:Per-Thread Statistical Counters}中
	那个显式的\co{counterp}数组不是又对线程数量
	施加了任意限制吗？
	为什么C语言不提供一个\co{per_thread()}接口，类似于
	Linux内核的\co{per_cpu()}原语，以便让线程
	更容易地访问其他线程的per-thread变量？
}\QuickQuizAnswer{
	确实如此，为什么不呢？

	公平地说，用户模式的thread-local存储面临一些
	Linux内核可以忽略的挑战。
	当用户级线程退出时，其per-thread变量全部消失，
	这使per-thread变量的访问变得复杂，
	特别是在用户级RCU出现之前
	（参见\cref{sec:defer:Read-Copy Update (RCU)}）。
	相比之下，在Linux内核中，当CPU离线时，
	该CPU的per-CPU变量仍然保持映射和可访问。

	类似地，当创建新的用户级线程时，
	其per-thread变量突然出现。
	相比之下，在Linux内核中，所有per-CPU变量
	在启动时就被映射和初始化，无论相应的CPU是否已经存在，
	甚至无论相应的CPU是否会存在。

	Linux内核施加的一个关键限制是CPU数量的编译时最大上限，
	即\co{CONFIG_NR_CPUS}，
	以及通常更紧的启动时上限\co{nr_cpu_ids}。
	相比之下，在用户空间中，线程数量不一定有硬编码的上限。

	当然，两种环境都必须处理动态加载的代码
	（用户空间的动态库，Linux内核的内核模块），
	这增加了per-thread变量的复杂性。

	这些复杂性使得用户空间环境提供对其他线程的
	per-thread变量的访问变得困难得多。
	尽管如此，这种访问非常有用，
	希望它有朝一日能出现。

	同时，像这样的教科书示例可以使用限制可由用户
	轻松调整的数组。
	或者，这样的数组可以在运行时动态分配和按需扩展。
	最后，可以使用可变长度数据结构（如链表），
	就像用户空间RCU库
	~\cite{MathieuDesnoyers2009URCU,MathieuDesnoyers2012URCU}中所做的那样。
	最后这种方法在某些情况下还可以减少\IX{false sharing}。
}\QuickQuizEnd

\begin{fcvref}[ln:count:count_end:whole:inc]
更新方调用的\co{inc_count()}函数非常简单，
如\clnrefrange{b}{e}所示。
\end{fcvref}

\begin{fcvref}[ln:count:count_end:whole:read]
读取方调用的\co{read_count()}函数稍微复杂一些。
\Clnref{acquire}获取lock以与正在退出的线程互斥，
\clnref{release}释放lock。
\Clnref{sum:init}初始化已退出线程的per-thread计数的总和，
\clnrefrange{loop:b}{loop:e}将还在运行的线程的per-thread计数累加进总和。
最后，\clnref{return}返回总和。
\end{fcvref}

\QuickQuizSeries{%
\QuickQuizB{
	\begin{fcvref}[ln:count:count_end:whole:read]
	\cref{lst:count:Per-Thread Statistical Counters}中
	\clnref{check}处对\co{NULL}的检查不会增加
	额外的分支预测错误吗？
	为什么不设一个永远为零的变量，让未使用的counter指针
	指向该变量而不是设为\co{NULL}？
	\end{fcvref}
}\QuickQuizAnswerB{
	这是一个合理的策略。
	检查性能差异留作读者的练习。
	但请记住，快速路径不是\co{read_count()}，
	而是\co{inc_count()}。
}\QuickQuizEndB
%
\QuickQuizE{
	为什么在
	\cref{lst:count:Per-Thread Statistical Counters}中
	的\co{read_count()}函数中需要像\emph{lock}这样
	重量级的东西来保护求和？
}\QuickQuizAnswerE{
	记住，当线程退出时，其per-thread变量会消失。
	因此，如果我们在线程退出后试图访问该线程的
	per-thread变量，将会得到段错误。
	lock协调求和和线程退出，防止这种情况发生。

	当然，我们可以使用读写lock来读获取，
	但\cref{chp:Deferred Processing}将介绍
	更轻量级的机制来实现所需的协调。

	另一种方法是使用数组而不是per-thread变量，
	正如Alexey Roytman指出的，这将消除
	对\co{NULL}的测试。
	然而，数组访问通常比per-thread变量的访问慢，
	而且使用数组意味着线程数量有固定上限。
	另外，请注意\co{inc_count()}快速路径上
	既不需要测试也不需要lock。
}\QuickQuizEndE
}

\begin{fcvref}[ln:count:count_end:whole:reg]
\Clnrefrange{b}{e}展示了\co{count_register_thread()}函数，
每个线程在首次使用此counter之前必须调用该函数。
该函数简单地将本线程对应\co{counterp[]}数组中的元素
指向线程的per-thread \co{counter}变量。
\end{fcvref}

\QuickQuiz{
	为什么在\cref{lst:count:Per-Thread Statistical Counters}中
	的\co{count_register_thread()}中需要获取lock？
	它只是对一个没有其他线程在修改的、正确对齐的
	机器字大小的位置进行存储，所以它应该本身就是
	atomic的，对吧？
}\QuickQuizAnswer{
	这个lock实际上可以省略，但安全起见更好，
	特别是考虑到这个函数仅在线程启动时执行，
	因此不在任何关键路径上。
	现在，如果我们在有数千个CPU的机器上测试，
	可能需要省略lock，但在``仅有''一百个左右CPU的机器上，
	没有必要花哨。
}\QuickQuizEnd

\begin{fcvref}[ln:count:count_end:whole:unreg]
\Clnrefrange{b}{e}展示了\co{count_unregister_thread()}函数，
之前调用过\co{count_register_thread()}的每个线程在退出前
必须调用该函数。
\Clnref{acquire}获取lock，
\clnref{release}释放它，从而排斥
\co{read_count()}的调用以及其他
\co{count_unregister_thread()}的调用。
\Clnref{add}将本线程的\co{counter}加到全局\co{finalcount}上，
然后\clnref{NULL}将\co{counterp[]}数组的对应元素设为\co{NULL}。
随后的\co{read_count()}调用可以在全局\co{finalcount}中找到
退出线程的计数值，并在遍历\co{counterp[]}数组时跳过已退出的线程，
这样就能获得正确的总数。
\end{fcvref}

这种方法让更新方的性能几乎和非atomic计数一样，并且也能线性扩展。
另一方面，并发的读取方争用一个全局lock，因此性能不佳，扩展能力也很差。
但是，对于statistical counters来说这不是问题，
因为statistical counters总是在增加计数，很少读取计数。
另外，由于给定线程的per-thread变量在该线程退出时消失，
本方法比基于数组的方案要复杂一些。

\QuickQuiz{
	好吧，但Linux内核在读取per-CPU counter的
	聚合值时不需要获取lock。
	那为什么用户空间代码需要这样做？
}\QuickQuizAnswer{
	记住，Linux内核的per-CPU变量始终是可访问的，
	即使相应的CPU离线——即使相应的CPU从未存在且永远不会存在。

\begin{listing}
\input{CodeSamples/count/count_tstat@whole.fcv}
\caption{带无锁（lock-free）求和的Per-Thread Statistical Counters}
\label{lst:count:Per-Thread Statistical Counters With Lockless Summation}
\end{listing}

	一种解决方法是确保每个线程继续存在直到所有线程完成，
	如\cref{lst:count:Per-Thread Statistical Counters With Lockless Summation}
	（\path{count_tstat.c}）所示。
	此代码的分析留作读者的练习，
	但请注意它需要对\path{counttorture.h} counter评估方案
	进行调整。
	（提示：参见\co{#ifndef KEEP_GCC_THREAD_LOCAL}。）
	\Cref{chp:Deferred Processing}将介绍
	以更加优雅的方式处理这种情况的同步机制。
}\QuickQuizEnd

上述方法都提供了出色的更新端性能和可扩展性。
然而，这种在更新端扩展极佳的方法，在存在大量线程时，会带来
读取端的巨大开销。
下一节将展示一种方法，能在保持更新端可扩展性的同时，
减少读取端产生的开销。

\subsection{最终一致性实现}
\label{sec:count:Eventually Consistent Implementation}

一种保持更新端可扩展性的同时又提升读取端性能的方法是
削弱一致性的要求。
前几节介绍的计数算法保证返回的值介于
\co{read_count()}执行前一刻的理想计数值和
执行完毕时的理想计数值之间。
\emph{最终一致性}~\cite{WernerVogels:2009:EventuallyConsistent}
提供了更弱的保证：
在没有\co{inc_count()}调用的情况下，
\co{read_count()}调用最终会返回准确的计数值。

我们通过维护一个全局counter来利用最终一致性。
但是因为写者只操纵它自己的per-thread counter，
所以单独有一个线程负责将per-thread counter的计数值传递给全局counter，
读者只是简单地访问全局counter的值。
如果写者正在更新计数，读者读出的值将不是最新的，
不过一旦写者更新完毕，全局counter最终会回归正确的值——
这就是为什么这种方法被称为最终一致性。

\begin{listing}
\ebresizeverb{1.0}{\input{CodeSamples/count/count_stat_eventual@whole.fcv}}
\caption{基于数组的Per-Thread最终一致性Counters}
\label{lst:count:Array-Based Per-Thread Eventually Consistent Counters}
\end{listing}

\begin{fcvref}[ln:count:count_stat_eventual:whole]
实现如
\cref{lst:count:Array-Based Per-Thread Eventually Consistent Counters}
（\path{count_stat_eventual.c}）所示。
\Clnrefrange{per_thr_cnt}{glb_cnt}
展示了跟踪counter值的per-thread变量和全局变量，
\clnref{stopflag}展示了用于协调终止的\co{stopflag}
（用于我们希望以准确counter值终止程序的情况）。
\clnrefrange{inc:b}{inc:e}所示的\co{inc_count()}函数类似于
\cref{lst:count:Array-Based Per-Thread Statistical Counters}中的对应函数。
\clnrefrange{read:b}{read:e}所示的\co{read_count()}函数
简单地返回\co{global_count}变量的值。

然而，\clnrefrange{init:b}{init:e}上的\co{count_init()}函数
创建了\clnrefrange{eventual:b}{eventual:e}所示的\co{eventual()}线程，
该线程循环遍历所有线程，
对per-thread本地\co{counter}求和并将结果存储到
\co{global_count}变量中。
\co{eventual()}线程在每次遍历之间等待任意选择的一毫秒。

\QuickQuiz{
	这种周期性扫描不会对电池寿命有害吗？
}\QuickQuizAnswer{
	很可能会。
	你能想到其他更节能的方式来更新\co{global_count}变量吗？
}\QuickQuizEnd

\clnrefrange{cleanup:b}{cleanup:e}上的\co{count_cleanup()}函数
协调终止。
这里的\co{smp_load_acquire()}调用和\co{eventual()}中的
\co{smp_store_release()}调用确保对\co{global_count}的所有更新
对\co{count_cleanup()}调用之后的代码可见。

本方法在提供极快的读取端性能的同时，仍然保持线性的
更新端可扩展性。但是，这种卓越的读取端性能和更新端可扩展性
带来的代价是需要一个额外的线程来运行\co{eventual()}。
\end{fcvref}

\QuickQuizSeries{%
\QuickQuizB{
	为什么\cref{lst:count:Array-Based Per-Thread Eventually Consistent Counters}中的
	\co{inc_count()}不需要使用atomic指令？
	毕竟，我们现在有多个线程访问per-thread counter！
}\QuickQuizAnswerB{
	因为两个线程中的一个只进行读取，并且由于
	变量是对齐的且是机器字大小的，非atomic指令即可满足。
	话虽如此，使用了\co{READ_ONCE()}宏来防止
	可能阻止counter更新对\co{eventual()}可见的
	编译器优化。\footnote{
		\co{READ_ONCE()}的一个简单定义见
		\cref{lst:toolsoftrade:Compiler Barrier Primitive (for GCC)}。}

	该算法的早期版本实际上使用了atomic指令，
	感谢Ersoy Bayramoglu指出它们实际上是不必要的。
	然而，请注意在32位系统上，
	per-thread \co{counter}变量可能需要限制为32位
	才能准确求和，但使用64位的\co{global_count}变量
	来避免溢出。
	在这种情况下，有必要周期性地将per-thread
	\co{counter}变量清零以避免溢出，
	这确实需要atomic指令。
	极其重要的是，这种清零不能延迟太久，
	否则较小的per-thread变量将会溢出。
	因此这种方法对底层系统施加了实时要求，
	必须极其谨慎地使用。

	相比之下，如果所有变量大小相同，
	任何变量的溢出都是无害的，因为最终的总和
	将是模字长的。
}\QuickQuizEndB
%
\QuickQuizM{
	\cref{lst:count:Array-Based Per-Thread Eventually Consistent Counters}中
	\co{eventual()}函数中的单个全局线程不会像全局lock一样
	成为严重的瓶颈吗？
}\QuickQuizAnswerM{
	在这种情况下不会。
	相反会发生的是，随着线程数量的增加，
	\co{read_count()}返回的counter值估计将变得更加不准确。
}\QuickQuizEndM
%
\QuickQuizM{
	\cref{lst:count:Array-Based Per-Thread Eventually Consistent Counters}中
	\co{read_count()}返回的估计值不会随着线程数量的增加
	而变得越来越不准确吗？
}\QuickQuizAnswerM{
	是的。
	如果这被证明是个问题，一种修复方法是提供多个
	\co{eventual()}线程，每个覆盖其他线程的一个子集。
	在更极端的情况下，可能需要\co{eventual()}线程的
	树状层次结构。
}\QuickQuizEndM
%
\QuickQuizM{
	鉴于在\cref{lst:count:Array-Based Per-Thread Eventually Consistent Counters}所示的
	最终一致性算法中，读取和更新的开销都极低且极具可扩展性，
	为什么有人还会使用
	\cref{sec:count:Array-Based Implementation}中描述的实现，
	考虑到它昂贵的读取端代码？
}\QuickQuizAnswerM{
	执行\co{eventual()}的线程消耗CPU时间。
	随着越来越多的最终一致性counter被添加，
	产生的\co{eventual()}线程最终将消耗所有可用CPU。
	因此，这种实现存在不同类型的可扩展性限制，
	该限制体现在最终一致性counter的数量上，
	而非线程或CPU的数量上。

	当然，可以做出其他权衡。
	例如，可以创建单个线程来处理所有最终一致性counter，
	这将把开销限制在单个CPU上，
	但随着counter数量的增加，更新到读取的延迟会增加。
	或者，该单个线程可以跟踪counter的更新速率，
	更频繁地访问频繁更新的counter。
	此外，处理counter的线程数量可以设置为
	CPU总数的某个比例，并且也许还可以在运行时调整。
	最后，每个counter可以指定其延迟，
	可以使用截止时间调度技术为每个counter
	提供所需的延迟。

	毫无疑问还有许多其他可能的权衡。
}\QuickQuizEndM
%
\QuickQuizE{
	\cref{lst:count:Array-Based Per-Thread Eventually Consistent Counters}中
	\co{read_count()}返回的估计值的准确度如何？
}\QuickQuizAnswerE{
	评估此估计的一种直接方法是使用
	\QuickQuizARef{\StatisticalCounterAccuracy}中推导的分析，
	但将$\Delta$设置为\co{eventual()}线程连续运行开始之间的间隔。
	处理给定counter具有多个\co{eventual()}线程的情况
	留作读者的练习。
}\QuickQuizEndE
}

\subsection{讨论}

上述实现表明，并发运行的statistical counters完全可以达到
接近单处理器的性能。

\QuickQuizSeries{%
\QuickQuizB{
	计数数据包和计数数据包中的总字节数之间有什么根本区别，
	假设数据包大小不同？
}\QuickQuizAnswerB{
	计数数据包时，counter仅递增值1。
	另一方面，计数字节时，counter可能递增
	相当大的数字。

	为什么这很重要？
	因为在递增值为1的情况下，返回的值在以下意义上是精确的：
	counter在某个时间点必然取过该值，
	即使无法准确说出该时间点何时发生。
	相比之下，在计数字节时，两个不同的线程可能返回
	与操作的任何全局排序不一致的值。

	为了看到这一点，假设线程~0向其counter加3，
	线程~1向其counter加5，
	线程~2和线程~3对counter求和。
	如果系统是``弱排序''的，或者编译器使用激进的优化，
	线程~2可能发现总和为3，线程~3可能发现总和为5。
	counter值序列仅有的全局顺序是0,3,8和0,5,8，
	两种顺序都与获得的结果不一致。

	如果你没答对这个问题，你并不孤单。
	Michael Scott在Paul E.~McKenney的博士答辩中
	用这个问题难住了Paul。
}\QuickQuizEndB
%
\QuickQuizE{
	鉴于读取方必须对所有线程的counter求和，
	当线程数量很大时，这个counter读取操作可能需要很长时间。
	有没有什么方法使递增操作保持快速和可扩展，
	同时让读取方也能享有不仅合理的性能和可扩展性，
	而且还有良好的准确性？
}\QuickQuizAnswerE{
	一种方法是维护一个全局近似值，
	类似于\cref{sec:count:Eventually Consistent Implementation}中描述的方法。
	更新方会递增其per-thread变量，但当它达到某个预定限制时，
	将其atomic地加到全局变量上，然后将per-thread变量清零。
	这将允许在平均递增开销和读取值的准确性之间进行权衡。
	特别是，它将允许对读取端不准确性设定明确的界限。

	另一种方法利用了这样一个事实：读取方通常只关心
	值的某些转变，而不是精确的值。
	这种方法在\cref{sec:count:Approximate Limit Counters}中进行了研究。

	鼓励读者想出并尝试其他方法，
	例如使用组合树。
}\QuickQuizEndE
}

通过学习本节的内容，你现在应该能回答
本章开头附近关于网络statistical counters的Quick Quiz了。

\section{近似限值计数器（limit counter）}
\label{sec:count:Approximate Limit Counters}
%
\epigraph{对正确问题的近似答案远比对近似问题的精确答案更有价值。}
	 {John Tukey}

计数的另一个特殊情况涉及上限检查。
例如，如\QuickQuizRef{\QcountQapproxcnt}中的
近似结构分配上限问题所述，
假设你需要维护一个已分配数据结构数目的计数器，
来防止分配超过一个上限（本例中为10,000）。
进一步假设这些结构的生命周期很短，数目极少能超出上限。
所以对近似上限来说，偶尔超出少许是可以接受的。
如果你需要上限必须精确，请参见\cref{sec:count:Exact Limit Counters}。

\subsection{设计}

实现limit counter的一种可能方法是将10,000的限制值平均划分
给每个线程，然后给每个线程一个固定个数的资源池。
假如有100个线程，每个线程管理一个有100个结构的资源池。
这种方法简单，在某些情况下也有效，但它无法处理
一种常见情况：某个结构由一个线程分配，但由另一个线程
释放~\cite{McKenney93}。
一方面，如果线程释放一个结构就得一分的话，
那么一直在分配的线程很快就耗尽了资源池，
而一直在释放的线程则积攒了大量分数却无法使用。
另一方面，如果每个被释放的结构都记在分配它的CPU名下，
CPU就需要操纵其他CPU的计数，这将会带来代价昂贵的
\IX{atomic}指令或其他跨线程通信手段。\footnote{
	话虽如此，如果每个结构总是由分配它的
	同一CPU（或线程）释放，那么这种简单的分区方法
	效果极好。}

简而言之，在很多重要的场景下我们都不能将counter完全分区。
在\cref{sec:count:Statistical Counters}中，考虑到对counter的分区才是
让更新端性能卓越的原因，这个事实可能会让我们感到悲观。
但是\cref{sec:count:Eventually Consistent Implementation}中
展示的最终一致性算法为我们提供了有趣的暗示。
该算法的关键是两套计数：
给更新方提供的per-thread \co{counter}变量和给读取方提供的
\co{global_count}变量，
还有一个定期更新\co{global_count}的\co{eventual()}线程，
使得结果最终与per-thread \co{counter}的值一致。
per-thread \co{counter}完美地分区了counter值，
而\co{global_count}保存完整的全局数值。

对于limit counters，我们也可以使用这种方法的变体，
部分地分区counter。
比如在四个线程中，每个线程都拥有一份per-thread \co{counter}，
但同时每个线程也持有一份per-thread的最大值\co{countermax}。

如果某个线程需要递增它的\co{counter}，
可是此时的\co{counter}等于\co{countermax}，该怎么办？
这里面有个技巧：此时可把此线程\co{counter}值的一半
转移给\co{globalcount}，然后再递增\co{counter}。
举个例子，假如某个线程的\co{counter}和\co{countermax}
现在都等于10，那么我们会执行如下操作：

\begin{enumerate}
\item	获取全局lock。
\item	给\co{globalcount}增加5。
\item	当前线程的\co{counter}减少5，以抵消全局的增加。
\item	释放全局lock。
\item	增加当前线程的\co{counter}，变成6。
\end{enumerate}

虽然这个办法仍然需要全局lock，但该lock只需要每五次递增操作
获取一次，这就大大降低了该lock的\IXalt{竞争}{lock contention}程度。
而且只要我们增大\co{countermax}的值，竞争程度还可以进一步降低。
不过，增大\co{countermax}值带来的代价是
\co{globalcount}精确度的降低。
为了说明这一点，假设我们有一台四CPU系统，此时\co{countermax}
的值为10，\co{globalcount}和真实计数值的误差最高可达40。
如果\co{countermax}的值增大到100，那么\co{globalcount}
和真实计数值的误差最高可达400。

这就提出了一个问题：我们到底有多在意\co{globalcount}
和真实计数值的偏差，其中真实计数值由\co{globalcount}
和所有per-thread \co{counter}变量相加得出。
这个问题的答案取决于真实计数值和计数上限
（这里我们将其命名为\co{globalcountmax}）的差值有多大。
这两个值的差值越大，\co{countermax}就越不容易
超过\co{globalcountmax}的上限。
这就意味着任何一个线程的\co{countermax}变量
可以根据当前的差值来设置。
当离上限还比较远时，可以给per-thread \co{countermax}赋一个较大的值，
这样对性能和可扩展性比较有好处；
当靠近上限时，可以给这些\co{countermax}赋一个较小的值，
以降低超过\co{globalcountmax}上限的风险。

这种设计就是一个\emph{parallel fastpath}的例子，
这是一种重要的设计模式，适用于以下情况：
在多数情况下没有线程间通信和交互的开销，
而偶尔进行的跨线程通信又使用了精心设计的
（但开销仍然很大的）全局算法。
此设计模式在\cref{sec:SMPdesign:Parallel Fastpath}中有更详细的介绍。

\subsection{简单限值计数器实现}
\label{sec:count:Simple Limit Counter Implementation}

\begin{fcvref}[ln:count:count_lim:variable]
\Cref{lst:count:Simple Limit Counter Variables}
展示了此实现使用的per-thread和全局变量。
per-thread \co{counter}和\co{countermax}变量分别是
相应线程的本地counter和该counter的上限。
\clnref{globalcountmax}上的\co{globalcountmax}变量
包含聚合counter的上限，
\clnref{globalcount}上的\co{globalcount}变量是全局counter。
\co{globalcount}和每个线程的\co{counter}之和给出了
整体counter的聚合值。
\clnref{globalreserve}上的\co{globalreserve}变量
至少是所有per-thread \co{countermax}变量之和。
\end{fcvref}
这些变量之间的关系如
\cref{fig:count:Simple Limit Counter Variable Relationships}所示：
\begin{enumerate}
\item	\co{globalcount}和\co{globalreserve}之和
	必须小于或等于\co{globalcountmax}。
\item	所有线程的\co{countermax}值之和
	必须小于或等于\co{globalreserve}。
\item	每个线程的\co{counter}必须小于或等于
	该线程的\co{countermax}。
\end{enumerate}

\begin{listing}
\input{CodeSamples/count/count_lim@variable.fcv}
\caption{简单Limit Counter变量}
\label{lst:count:Simple Limit Counter Variables}
\end{listing}

\begin{figure}
\centering
\resizebox{2.5in}{!}{\includegraphics{count/count_lim}}
\caption{简单Limit Counter变量关系}
\label{fig:count:Simple Limit Counter Variable Relationships}
\end{figure}

\co{counterp[]}数组的每个元素指向对应线程的\co{counter}变量。
最后，\co{gblcnt_mutex} spinlock保护所有全局变量，
也就是说，除非线程获取了\co{gblcnt_mutex}，
否则不能访问或修改任何全局变量。

\begin{listing}
\input{CodeSamples/count/count_lim@add_sub_read.fcv}
\caption{简单Limit Counter的加法、减法和读取}
\label{lst:count:Simple Limit Counter Add; Subtract; and Read}
\end{listing}

\Cref{lst:count:Simple Limit Counter Add; Subtract; and Read}
展示了\co{add_count()}、\co{sub_count()}和\co{read_count()}
函数（\path{count_lim.c}）。

\QuickQuiz{
	为什么\cref{lst:count:Simple Limit Counter Add; Subtract; and Read}
	提供\co{add_count()}和\co{sub_count()}而不是
	\cref{sec:count:Statistical Counters}中所示的
	\co{inc_count()}和\co{dec_count()}接口？
}\QuickQuizAnswer{
	因为结构有不同的大小。
	当然，对应于特定大小结构的limit counter
	可能仍然可以使用\co{inc_count()}和\co{dec_count()}。
}\QuickQuizEnd

\begin{fcvref}[ln:count:count_lim:add_sub_read:add]
\Clnrefrange{b}{e}展示了\co{add_count()}，
它将指定值\co{delta}加到counter上。
\Clnref{checklocal}检查此线程的\co{counter}上
是否有空间容纳\co{delta}，
如果有，\clnref{add}执行加法，\clnref{return:ls}返回成功。
这是\co{add_counter()}的快速路径，它不执行任何atomic操作，
仅引用per-thread变量，不应引起任何cache miss。
\end{fcvref}

\begin{listing}
\begin{VerbatimL}[firstnumber=3]
	if (counter + delta <= countermax) {
		WRITE_ONCE(counter, counter + delta);
		return 1;
	}
\end{VerbatimL}
\caption{直觉上的快速路径}
\label{lst:count:Intuitive Fastpath}
\end{listing}

\QuickQuiz{
	\cref{lst:count:Simple Limit Counter Add; Subtract; and Read}中
	\clnrefr{ln:count:count_lim:add_sub_read:add:checklocal}
	上那个奇怪形式的条件是怎么回事？
	为什么不使用\cref{lst:count:Intuitive Fastpath}中所示的
	更直观形式的快速路径？
}\QuickQuizAnswer{
	两个词：``整数溢出。''

	用\co{counter}等于~10、
	\co{delta}等于\co{ULONG_MAX}来试试
	\cref{lst:count:Intuitive Fastpath}中的公式。
	然后再用
	\cref{lst:count:Simple Limit Counter Add; Subtract; and Read}
	中所示的代码试一试。

	要理解本示例的其余部分，需要对整数溢出有良好的理解，
	因此如果你以前从未处理过整数溢出，
	请尝试几个例子来掌握要领。
	整数溢出有时可能比并行算法更难处理正确！
}\QuickQuizEnd

\begin{fcvref}[ln:count:count_lim:add_sub_read:add]
如果\clnref{checklocal}上的测试失败，我们必须访问全局变量，
因此必须在\clnref{acquire}上获取\co{gblcnt_mutex}，
在失败情况下于\clnref{release:f}释放，
在成功情况下于\clnref{release:s}释放。
\Clnref{globalize}调用\co{globalize_count()}
（见\cref{lst:count:Simple Limit Counter Utility Functions}），
该函数清零线程本地变量，并根据需要调整全局变量，
这就简化了全局处理过程。
（别太相信我的话，试着自己写代码！）
\Clnref{checkglb:b,checkglb:e}检查是否能容纳\co{delta}的加法，
小于号前面表达式的含义如
\cref{fig:count:Simple Limit Counter Variable Relationships}中
两个红色（最左边的）柱子的高度差所示。
如果无法容纳\co{delta}的加法，
则\clnref{release:f}（如前所述）释放\co{gblcnt_mutex}，
\clnref{return:gf}返回失败。

否则，我们走慢速路径。
\Clnref{addglb}将\co{delta}加到\co{globalcount}上，然后
\clnref{balance}调用\co{balance_count()}
（见\cref{lst:count:Simple Limit Counter Utility Functions}）
以更新全局和per-thread变量。
这次\co{balance_count()}调用通常会设置此线程的
\co{countermax}以重新启用快速路径。
\Clnref{release:s}然后释放\co{gblcnt_mutex}（同样如前所述），
最后\clnref{return:gs}返回成功。
\end{fcvref}

\QuickQuiz{
	为什么在\cref{lst:count:Simple Limit Counter Add; Subtract; and Read}中
	\co{globalize_count()}先将per-thread变量清零，
	然后再调用\co{balance_count()}来重新填充它们？
	为什么不直接让per-thread变量保持非零？
}\QuickQuizAnswer{
	实际上，该代码的早期版本就是这么做的。
	但加法和减法极其廉价，而处理所有出现的特殊情况
	相当复杂。
	同样，请随意自己尝试，但要小心整数溢出！
}\QuickQuizEnd

\begin{fcvref}[ln:count:count_lim:add_sub_read:sub]
\Clnrefrange{b}{e}展示了\co{sub_count()}，
它从counter中减去指定的\co{delta}。
\Clnref{checklocal}检查per-thread counter是否能容纳此次减法，
如果能，\clnref{sub}执行减法，
\clnref{return:ls}返回成功。
这些行构成了\co{sub_count()}的快速路径，
与\co{add_count()}一样，此快速路径不执行任何昂贵操作。

如果快速路径无法容纳\co{delta}的减法，
执行继续到\clnrefrange{acquire}{return:gs}的慢速路径。
因为慢速路径必须访问全局状态，\clnref{acquire}
获取\co{gblcnt_mutex}，它由\clnref{release:f}
（失败情况）或\clnref{release:s}（成功情况）释放。
\Clnref{globalize}调用\co{globalize_count()}
（见\cref{lst:count:Simple Limit Counter Utility Functions}），
同样清除线程本地变量并根据需要调整全局变量。
\Clnref{checkglb}检查counter是否能容纳\co{delta}的减法，
如果不能，\clnref{release:f}释放\co{gblcnt_mutex}
（如前所述），\clnref{return:gf}返回失败。
\end{fcvref}

\QuickQuizSeries{%
\QuickQuizB{
	既然\co{globalreserve}在\co{add_count()}中对我们不利，
	为什么它在
	\cref{lst:count:Simple Limit Counter Add; Subtract; and Read}
	的\co{sub_count()}中不对我们有利呢？
}\QuickQuizAnswerB{
	\co{globalreserve}变量跟踪所有线程\co{countermax}变量之和。
	这些线程的\co{counter}变量之和可能在零到
	\co{globalreserve}之间的任何位置。
	因此我们必须采取保守的方法，
	在\co{add_count()}中假设所有线程的\co{counter}变量都是满的，
	在\co{sub_count()}中假设它们都是空的。

	但请记住这个问题，我们稍后会回来讨论它。
}\QuickQuizEndB
%
\QuickQuizE{
	假设一个线程调用了
	\cref{lst:count:Simple Limit Counter Add; Subtract; and Read}
	中的\co{add_count()}，然后另一个线程调用\co{sub_count()}。
	\co{sub_count()}不会返回失败，即使counter的值不为零？
}\QuickQuizAnswerE{
	确实会！
	在许多情况下，这将是一个问题，
	如\cref{sec:count:Simple Limit Counter Discussion}所讨论的，
	在这些情况下，\cref{sec:count:Exact Limit Counters}
	中的算法可能更可取。
}\QuickQuizEndE
}

\begin{fcvref}[ln:count:count_lim:add_sub_read:sub]
另一方面，如果\clnref{checkglb}发现counter\emph{能}容纳
\co{delta}的减法，我们完成慢速路径。
\Clnref{subglb}执行减法，然后
\clnref{balance}调用\co{balance_count()}
（见\cref{lst:count:Simple Limit Counter Utility Functions}）
以更新全局和per-thread变量
（希望能重新启用快速路径）。
然后\clnref{release:s}释放\co{gblcnt_mutex}，
\clnref{return:gs}返回成功。
\end{fcvref}

\QuickQuiz{
	为什么在\cref{lst:count:Simple Limit Counter Add; Subtract; and Read}中
	同时有\co{add_count()}和\co{sub_count()}？
	为什么不简单地向\co{add_count()}传递一个负数？
}\QuickQuizAnswer{
	鉴于\co{add_count()}接受\co{unsigned} \co{long}
	作为参数，向它传递负数将会相当困难。
	而且除非你有什么反物质内存，
	在计算使用中的结构数量时允许负数也没什么意义！

	撇开玩笑不谈，当然可以合并
	\co{add_count()}和\co{sub_count()}，但合并函数中的
	\co{if}条件将比当前这对函数中的更复杂，
	这反过来意味着这些快速路径的执行会更慢。
}\QuickQuizEnd

\begin{fcvref}[ln:count:count_lim:add_sub_read:read]
\Clnrefrange{b}{e}展示了\co{read_count()}，
它返回counter的聚合值。
它在\clnref{acquire}获取\co{gblcnt_mutex}，
在\clnref{release}释放它，
排斥\co{add_count()}和\co{sub_count()}的全局操作，
并且如我们将看到的，还排斥线程创建和退出。
\Clnref{initsum}将局部变量\co{sum}初始化为\co{globalcount}的值，
然后\clnrefrange{loop:b}{loop:e}的循环
对per-thread \co{counter}变量求和。
\Clnref{return}然后返回总和。
\end{fcvref}

\begin{listing}
\input{CodeSamples/count/count_lim@utility.fcv}
\caption{简单Limit Counter工具函数}
\label{lst:count:Simple Limit Counter Utility Functions}
\end{listing}

\Cref{lst:count:Simple Limit Counter Utility Functions}
展示了\co{add_count()}、\co{sub_count()}和\co{read_count()}
原语（见\cref{lst:count:Simple Limit Counter Add; Subtract; and Read}）
使用的一些工具函数。

\begin{fcvref}[ln:count:count_lim:utility:globalize]
\Clnrefrange{b}{e}展示了\co{globalize_count()}，
它将当前线程的per-thread counter清零，并适当调整全局变量。
重要的是，此函数不改变counter的聚合值，
而是改变counter当前值的表示方式。
\Clnref{add}将线程的\co{counter}变量加到\co{globalcount}上，
\clnref{zero}将\co{counter}清零。
类似地，\clnref{sub}从\co{globalreserve}中减去per-thread
\co{countermax}，\clnref{zeromax}将\co{countermax}清零。
阅读此函数和下面的\co{balance_count()}时，
参考\cref{fig:count:Simple Limit Counter Variable Relationships}
会有所帮助。
\end{fcvref}

\begin{fcvref}[ln:count:count_lim:utility:balance]
\Clnrefrange{b}{e}展示了\co{balance_count()}，
粗略地说是\co{globalize_count()}的逆操作。
此函数的任务是将当前线程的\co{countermax}变量
设置为最大值，同时避免counter超过\co{globalcountmax}限制的风险。
改变当前线程的\co{countermax}变量当然需要
对\co{counter}、\co{globalcount}和\co{globalreserve}
做相应调整，如参考
\cref{fig:count:Simple Limit Counter Variable Relationships}
所示。
通过这样做，\co{balance_count()}最大化了
\co{add_count()}和\co{sub_count()}低开销快速路径的使用。
与\co{globalize_count()}一样，\co{balance_count()}
不允许改变counter的聚合值。

\Clnrefrange{share:b}{share:e}计算此线程在
\co{globalcountmax}中尚未被\co{globalcount}或\co{globalreserve}
覆盖的部分的份额，并将计算出的数量赋给此线程的\co{countermax}。
\Clnref{adjreserve}对\co{globalreserve}做相应调整。
\Clnref{middle}将此线程的\co{counter}设置为零到\co{countermax}
范围的中间值。
\Clnref{check}检查\co{globalcount}是否实际上能容纳
此\co{counter}值，如果不能，
\clnref{adjcounter}相应减少\co{counter}。
最后，在两种情况下，\clnref{adjglobal}对\co{globalcount}
做相应调整。
\end{fcvref}

\QuickQuiz{
	\begin{fcvref}[ln:count:count_lim:utility:balance]
	为什么在\cref{lst:count:Simple Limit Counter Utility Functions}的
	\clnref{middle}将\co{counter}设置为\co{countermax / 2}？
	直接取\co{countermax}次计数不是更简单吗？
	\end{fcvref}
}\QuickQuizAnswer{
	\begin{fcvref}[ln:count:count_lim:utility:balance]
	首先，它确实保留了\co{countermax}次计数
	（见\clnref{adjreserve}），但是
	它调整为当前线程实际只使用其中一半。
	这允许线程执行至少\co{countermax / 2}次
	递增或递减，然后才需要再次引用\co{globalcount}。

	注意\co{globalcount}中的记账保持准确，
	这要归功于\clnref{adjglobal}的调整。
	\end{fcvref}
}\QuickQuizEnd

\begin{figure*}
\centering
\IfEbookSize{
\resizebox{\textwidth}{!}{\includegraphics{count/globbal}}
}{
\resizebox{5in}{!}{\includegraphics{count/globbal}}
}
\caption{全局化和平衡的示意图}
\label{fig:count:Schematic of Globalization and Balancing}
\end{figure*}

看一个描绘counter关系如何随
\co{globalize_count()}和\co{balance_count()}的执行
而变化的示意图是很有帮助的，
如\cref{fig:count:Schematic of Globalization and Balancing}所示。
时间从左到右推进，最左边的配置大致是
\cref{fig:count:Simple Limit Counter Variable Relationships}的配置。
中间的配置展示了线程~0执行\co{globalize_count()}后
这些counter之间的关系。
从图中可以看出，线程~0的\co{counter}（图中的``c~0''）
被加到\co{globalcount}上，而\co{globalreserve}的值
减少了相同的量。
线程~0的\co{counter}和\co{countermax}
（图中的``cm~0''）都减少到零。
其他三个线程的counter不变。
注意此更改没有影响counter的整体值，
如连接最左和中间配置的最下方虚线所示。
换句话说，\co{globalcount}和四个线程的\co{counter}变量之和
在两个配置中相同。
类似地，此更改没有影响\co{globalcount}和\co{globalreserve}之和，
如上方虚线所示。

最右边的配置展示了线程~0再次执行\co{balance_count()}后
这些counter之间的关系。
剩余计数的四分之一（由从所有三个配置向上延伸的垂直线表示）
被加到线程~0的\co{countermax}上，其中一半加到线程~0的\co{counter}上。
加到线程~0的\co{counter}上的量也从\co{globalcount}中减去，
以避免改变counter的整体值
（它仍然是\co{globalcount}和三个线程的\co{counter}变量之和），
再次如连接中间和最右配置的两条虚线中较低的一条所示。
\co{globalreserve}变量也被调整，使此变量保持等于
四个线程的\co{countermax}变量之和。
因为线程~0的\co{counter}小于其\co{countermax}，
线程~0可以再次本地递增counter。

\QuickQuiz{
	在\cref{fig:count:Schematic of Globalization and Balancing}中，
	尽管剩余计数到限制的四分之一被分配给线程~0，
	但只有八分之一的剩余计数被消耗了，
	如连接中间和最右配置的最上方虚线所示。
	为什么会这样？
}\QuickQuizAnswer{
	这发生的原因是线程~0的\co{counter}被设置为
	其\co{countermax}的一半。
	因此，在分配给线程~0的四分之一中，
	该四分之一的一半（即八分之一）来自\co{globalcount}，
	另一半（同样是八分之一）来自剩余计数。

	采用这种方法有两个目的：
	\begin{enumerate*}[(1)]
	\item 允许线程~0对递减和递增都使用快速路径，
	\item 如果所有线程都在单调递增趋向限制时减少不准确性。
	\end{enumerate*}
	要看到最后一点，请逐步执行算法并观察它的行为。
}\QuickQuizEnd

\begin{fcvref}[ln:count:count_lim:utility:register]
\Clnrefrange{b}{e}展示了\co{count_register_thread()}，
它为新创建的线程设置状态。
此函数简单地在\co{gblcnt_mutex}保护下
将指向新创建线程的\co{counter}变量的指针
安装到\co{counterp[]}数组的相应条目中。
\end{fcvref}

\begin{fcvref}[ln:count:count_lim:utility:unregister]
最后，\clnrefrange{b}{e}展示了\co{count_unregister_thread()}，
它为即将退出的线程拆除状态。
\Clnref{acquire}获取\co{gblcnt_mutex}，
\clnref{release}释放它。
\Clnref{globalize}调用\co{globalize_count()}
清除此线程的counter状态，
\clnref{clear}清除此线程在\co{counterp[]}数组中的条目。
\end{fcvref}

\subsection{简单限值计数器讨论}
\label{sec:count:Simple Limit Counter Discussion}

当聚合值接近零时，这种counter运行得相当快，
只在\co{add_count()}和\co{sub_count()}快速路径中的
比较和分支上存在一些开销。
但是，per-thread \co{countermax}的使用表明
即使counter的聚合值离\co{globalcountmax}很远，
\co{add_count()}也可能失败。
同样，即使counter的聚合值远大于零，
\co{sub_count()}也可能失败。

在许多情况下，这都是不可接受的。
即使\co{globalcountmax}只不过是一个近似的上限，
一般也有一个近似度的容忍限度。
一种限制近似度的方法是对per-thread \co{countermax}的值
强加一个上限。
这个任务在下一节中完成。

\subsection{近似限值计数器实现}
\label{sec:count:Approximate Limit Counter Implementation}

因为此实现（\path{count_lim_app.c}）与前一节中的实现
（\cref{lst:count:Simple Limit Counter Variables,%
lst:count:Simple Limit Counter Add; Subtract; and Read,%
lst:count:Simple Limit Counter Utility Functions}）
非常相似，这里只展示差异。
\Cref{lst:count:Approximate Limit Counter Variables}
与\cref{lst:count:Simple Limit Counter Variables}相同，
增加了\co{MAX_COUNTERMAX}，它设置per-thread
\co{countermax}变量的最大允许值。

\begin{listing}
\input{CodeSamples/count/count_lim_app@variable.fcv}
\caption{Approximate Limit Counter变量}
\label{lst:count:Approximate Limit Counter Variables}
\end{listing}

\begin{listing}
\input{CodeSamples/count/count_lim_app@balance.fcv}
\caption{Approximate Limit Counter平衡}
\label{lst:count:Approximate Limit Counter Balancing}
\end{listing}

\begin{fcvref}[ln:count:count_lim_app:balance]
类似地，\cref{lst:count:Approximate Limit Counter Balancing}
与\cref{lst:count:Simple Limit Counter Utility Functions}中的
\co{balance_count()}函数相同，
增加了\clnref{enforce:b,enforce:e}，
它们对per-thread \co{countermax}变量强制执行
\co{MAX_COUNTERMAX}限制。
\end{fcvref}

\subsection{近似限值计数器讨论}

上述改动极大地减小了前一版本中出现的上限不准确程度，
但又带来了另一个问题：
任何给定大小的\co{MAX_COUNTERMAX}将导致取决于工作负载的
一部分访问无法进入快速路径。
随着线程数的增加，非快速路径的执行将成为
性能和可扩展性上的问题。
不过，我们暂时放下这个问题，先去看看具有精确上限的counter。

\section{精确限值计数器}
\label{sec:count:Exact Limit Counters}
%
\epigraph{精确是有代价的。请明智地花费。}{佚名}

要解决\QuickQuizRef{\QcountQexactcnt}中提出的
精确结构分配上限问题，
我们需要一个能精确判断何时计数超过上限的limit counter。
一种实现的办法是允许线程放弃自己的保留计数，
另一种做法是采用\IX{atomic}指令。
当然，atomic指令会减慢快速路径，但另一方面，
连试都不试又有点过于愚蠢了。

\subsection{原子限值计数器实现}
\label{sec:count:Atomic Limit Counter Implementation}

不幸的是，如果想要某个线程减少另一个线程上的计数，
需要atomic地操纵两个线程的\co{counter}和\co{countermax}变量。
通常的做法是将这两个变量合并为单个变量，
比如，一个32位变量，高16位代表\co{counter}，
低16位代表\co{countermax}。

\QuickQuiz{
	为什么需要将线程的\co{counter}和\co{countermax}变量
	作为一个整体进行atomic操作？
	单独对它们进行atomic操作不够吗？
}\QuickQuizAnswer{
	这也许是可能的，但需要非常小心。
	注意，如果不先将\co{countermax}清零就移除\co{counter}，
	可能导致相应线程在\co{counter}被清零后立即增加它，
	完全抵消了清零counter的效果。

	相反的顺序，即先清零\co{countermax}再移除\co{counter}，
	也可能导致\co{counter}非零。
	要看到这一点，考虑以下事件序列：

	\begin{enumerate}
	\item	线程~A获取其\co{countermax}，发现它不为零。
	\item	线程~B将线程~A的\co{countermax}清零。
	\item	线程~B移除线程~A的\co{counter}。
	\item	线程~A发现其\co{countermax}不为零后，
		继续向其\co{counter}加值，
		导致\co{counter}值非零。
	\end{enumerate}

	同样，分别对\co{countermax}和\co{counter}进行
	atomic操作也许是可能的，
	但显然需要非常小心。
	而且这样做很可能会减慢快速路径。

	探索这些可能性留作读者的练习。
}\QuickQuizEnd

\begin{listing}
\input{CodeSamples/count/count_lim_atomic@var_access.fcv}
\caption{Atomic Limit Counter变量和访问函数}
\label{lst:count:Atomic Limit Counter Variables and Access Functions}
\end{listing}

\begin{fcvref}[ln:count:count_lim_atomic:var_access:var]
一个简单的atomic limit counter的变量和访问函数如
\cref{lst:count:Atomic Limit Counter Variables and Access Functions}
（\path{count_lim_atomic.c}）所示。
早期算法中的\co{counter}和\co{countermax}变量
被合并为\clnref{candmax}上的单个变量\co{counterandmax}，
\co{counter}在高半部分，\co{countermax}在低半部分。
此变量的类型为\apik{atomic_t}，其底层表示为\co{int}。

\Clnrefrange{def:b}{def:e}展示了\co{globalcountmax}、\co{globalcount}、
\co{globalreserve}、\co{counterp}和\co{gblcnt_mutex}的定义，
它们的角色与
\cref{lst:count:Approximate Limit Counter Variables}中的对应部分类似。
\Clnref{CM_BITS}定义了\co{CM_BITS}，它给出\co{counterandmax}
每半部分的位数，\clnref{MAX_CMAX}定义了\co{MAX_COUNTERMAX}，
它给出\co{counterandmax}任一半部分可容纳的最大值。
\end{fcvref}

\QuickQuiz{
	\cref{lst:count:Atomic Limit Counter Variables and Access Functions}的
	\clnrefr{ln:count:count_lim_atomic:var_access:var:CM_BITS}
	在什么方面违反了C标准？
}\QuickQuizAnswer{
	它假设每字节八位。
	这个假设对所有当前可以轻松组装成共享内存多处理器的
	商用微处理器都成立，但对所有曾经运行过C代码的
	计算机系统并非如此。
	（你可以做什么来替代以符合C标准？会有什么缺点？）
}\QuickQuizEnd

\begin{fcvref}[ln:count:count_lim_atomic:var_access:split_int]
\Clnrefrange{b}{e}展示了\co{split_counterandmax_int()}函数，
当给定\co{atomic_t counterandmax}变量的底层\co{int}时，
将其拆分为\co{counter}（\co{c}）
和\co{countermax}（\co{cm}）组件。
\Clnref{msh}隔离此\co{int}的最高有效半部分，
将结果放在参数\co{c}指定的位置，
\clnref{lsh}隔离此\co{int}的最低有效半部分，
将结果放在参数\co{cm}指定的位置。
\end{fcvref}

\begin{fcvref}[ln:count:count_lim_atomic:var_access:split]
\Clnrefrange{b}{e}展示了\co{split_counterandmax()}函数，
它在\clnref{int}获取指定变量的底层\co{int}，
在\clnref{old}将其存储为\co{old}参数所指定的位置，
然后在\clnref{split_int}调用\co{split_counterandmax_int()}
进行拆分。
\end{fcvref}

\QuickQuiz{
	既然只有一个\co{counterandmax}变量，
	为什么在\cref{lst:count:Atomic Limit Counter Variables and Access Functions}的
	\clnrefr{ln:count:count_lim_atomic:var_access:split:func}上
	还要传入指向它的指针？
}\QuickQuizAnswer{
	每个线程有一个\co{counterandmax}变量。
	稍后，我们将看到需要将其他线程的
	\co{counterandmax}变量传递给\co{split_counterandmax()}的代码。
}\QuickQuizEnd

\begin{fcvref}[ln:count:count_lim_atomic:var_access:merge]
\Clnrefrange{b}{e}展示了\co{merge_counterandmax()}函数，
它可以被看作\co{split_counterandmax()}的逆操作。
\Clnref{merge}合并通过\co{c}和\co{cm}传入的
\co{counter}和\co{countermax}值，并返回结果。
\end{fcvref}

\QuickQuiz{
	为什么\cref{lst:count:Atomic Limit Counter Variables and Access Functions}中的
	\co{merge_counterandmax()}返回\co{int}
	而不是直接存储到\co{atomic_t}中？
}\QuickQuizAnswer{
	稍后，我们将看到需要将\co{int}返回值
	传递给\co{atomic_cmpxchg()}原语。
}\QuickQuizEnd

\begin{listing}
\ebresizeverb{.81}{\input{CodeSamples/count/count_lim_atomic@add_sub.fcv}}
\caption{Atomic Limit Counter加法和减法}
\label{lst:count:Atomic Limit Counter Add and Subtract}
\end{listing}

\Cref{lst:count:Atomic Limit Counter Add and Subtract}
展示了\co{add_count()}和\co{sub_count()}函数。

\begin{fcvref}[ln:count:count_lim_atomic:add_sub:add]
\Clnrefrange{b}{e}展示了\co{add_count()}，其快速路径覆盖
\clnrefrange{fast:b}{return:fs}，
函数的其余部分为慢速路径。
快速路径的\clnrefrange{fast:b}{loop:e}构成一个比较并交换
（CAS）循环，
\clnrefrange{atmcmpex}{loop:e}上的\apik{atomic_cmpxchg()}原语
执行实际的CAS。
\Clnref{split}将当前线程的\co{counterandmax}变量拆分为
\co{counter}（\co{c}中）和\co{countermax}（\co{cm}中）组件，
同时将底层\co{int}放入\co{old}中。
\Clnref{check}检查\co{delta}量是否能在本地容纳
（注意避免整数溢出），如果不能，
\clnref{goto}转到慢速路径。
否则，\clnref{merge}将更新后的\co{counter}值与
原始\co{countermax}值合并为\co{new}。
\clnrefrange{atmcmpex}{loop:e}上的\co{atomic_cmpxchg()}原语
然后atomic地将此线程的\co{counterandmax}变量与\co{old}比较，
如果比较成功则将其值更新为\co{new}。
如果比较成功，\clnref{return:fs}返回成功，
否则执行在\clnref{fast:b}继续循环。
\end{fcvref}

\QuickQuizSeries{%
\QuickQuizB{
	呸！
	\cref{lst:count:Atomic Limit Counter Add and Subtract}
	\clnrefr{ln:count:count_lim_atomic:add_sub:add:goto}上
	那个丑陋的\co{goto}是怎么回事？
	你没听说过\co{break}语句吗？
}\QuickQuizAnswerB{
	用\co{break}替换\co{goto}需要保留一个标志
	来确定是否应该在
	\clnrefr{ln:count:count_lim_atomic:add_sub:add:return:fs}
	处返回，
	这不是你在快速路径上想要的东西。
	如果你真的那么讨厌\co{goto}，最好的办法是
	将快速路径提取到一个返回成功或失败的单独函数中，
	``失败''表示需要慢速路径。
	这留作讨厌goto的读者的练习。
}\QuickQuizEndB
%
\QuickQuizE{
	\begin{fcvref}[ln:count:count_lim_atomic:add_sub:add]
	\cref{lst:count:Atomic Limit Counter Add and Subtract}
	\clnrefrange{atmcmpex}{loop:e}上的
	\co{atomic_cmpxchg()}原语为什么会失败？
	毕竟，我们在\clnref{split}上获取了它的旧值，
	而且没有改变它！
	\end{fcvref}
}\QuickQuizAnswerE{
	\begin{fcvref}[ln:count:count_lim_atomic:add_sub:add]
	稍后，我们将看到
	\cref{lst:count:Atomic Limit Counter Utility Functions 1}中的
	\co{flush_local_count()}函数如何可能与
	\cref{lst:count:Atomic Limit Counter Add and Subtract}
	\clnrefrange{fast:b}{loop:e}上的快速路径
	并发更新此线程的\co{counterandmax}变量。
	\end{fcvref}
}\QuickQuizEndE
}

\begin{fcvref}[ln:count:count_lim_atomic:add_sub:add]
\cref{lst:count:Atomic Limit Counter Add and Subtract}的
\Clnrefrange{slow:b}{return:ss}
展示了\co{add_count()}的慢速路径，
它由\co{gblcnt_mutex}保护，
在\clnref{acquire}获取，在\clnref{release:f,release:s}释放。
\Clnref{globalize}调用\co{globalize_count()}，
将此线程的状态移到全局counter。
\Clnrefrange{checkglb:b}{checkglb:e}检查当前全局状态
是否能容纳\co{delta}值，如果不能，\clnref{flush}调用
\co{flush_local_count()}将所有线程的本地状态刷新到
全局counter，然后\clnrefrange{checkglb:nb}{checkglb:ne}
重新检查是否能容纳\co{delta}。
如果经过所有这些之后，\co{delta}的加法仍然无法容纳，
则\clnref{release:f}释放\co{gblcnt_mutex}（如前所述），
然后\clnref{return:sf}返回失败。

否则，\clnref{addglb}将\co{delta}加到全局counter上，
\clnref{balance}在适当时将计数分配到本地状态，
\clnref{release:s}释放\co{gblcnt_mutex}（同样如前所述），
最后\clnref{return:ss}返回成功。
\end{fcvref}

\begin{fcvref}[ln:count:count_lim_atomic:add_sub:sub]
\cref{lst:count:Atomic Limit Counter Add and Subtract}的
\Clnrefrange{b}{e}
展示了\co{sub_count()}，其结构与\co{add_count()}类似，
快速路径在\clnrefrange{fast:b}{fast:e}，
慢速路径在\clnrefrange{slow:b}{slow:e}。
此函数的逐行分析留作读者的练习。
\end{fcvref}

\begin{listing}
\input{CodeSamples/count/count_lim_atomic@read.fcv}
\caption{Atomic Limit Counter读取}
\label{lst:count:Atomic Limit Counter Read}
\end{listing}

\begin{fcvref}[ln:count:count_lim_atomic:read]
\Cref{lst:count:Atomic Limit Counter Read}展示了\co{read_count()}。
\Clnref{acquire}获取\co{gblcnt_mutex}，
\clnref{release}释放它。
\Clnref{initsum}将局部变量\co{sum}初始化为\co{globalcount}的值，
\clnrefrange{loop:b}{loop:e}的循环将per-thread counter
加到此总和中，在\clnref{split}使用\co{split_counterandmax}
隔离每个per-thread counter。
最后，\clnref{return}返回总和。
\end{fcvref}

\begin{listing}
\input{CodeSamples/count/count_lim_atomic@utility1.fcv}
\caption{Atomic Limit Counter工具函数1}
\label{lst:count:Atomic Limit Counter Utility Functions 1}
\end{listing}

\begin{listing}
\input{CodeSamples/count/count_lim_atomic@utility2.fcv}
\caption{Atomic Limit Counter工具函数2}
\label{lst:count:Atomic Limit Counter Utility Functions 2}
\end{listing}

\Cref{lst:count:Atomic Limit Counter Utility Functions 1,%
lst:count:Atomic Limit Counter Utility Functions 2}
展示了工具函数
\co{globalize_count()}、
\co{flush_local_count()}、
\co{balance_count()}、
\co{count_register_thread()}和
\co{count_unregister_thread()}。
\begin{fcvref}[ln:count:count_lim_atomic:utility1:globalize]
\co{globalize_count()}的代码在
\cref{lst:count:Atomic Limit Counter Utility Functions 1}的
\clnrefrange{b}{e}上展示，
与前面算法类似，增加了\clnref{split}，
现在需要它来从\co{counterandmax}中拆分出
\co{counter}和\co{countermax}。
\end{fcvref}

\begin{fcvref}[ln:count:count_lim_atomic:utility1:flush]
\co{flush_local_count()}的代码在
\clnrefrange{b}{e}上展示，
它将所有线程的本地counter状态移到全局counter。
\Clnref{checkrsv}检查\co{globalreserve}的值
是否允许任何per-thread计数，如果不允许，\clnref{return:n}返回。
否则，\clnref{initzero}将局部变量\co{zero}初始化为
合并的零\co{counter}和\co{countermax}。
\clnrefrange{loop:b}{loop:e}的循环遍历每个线程。
\Clnref{checkp}检查当前线程是否有counter状态，
如果有，\clnrefrange{atmxchg}{glbrsv}将该状态
移到全局counter。
\Clnref{atmxchg}在用零替换的同时atomic地获取
当前线程的状态。
\Clnref{split}将此状态拆分为\co{counter}
（局部变量\co{c}中）和\co{countermax}（局部变量\co{cm}中）组件。
\Clnref{glbcnt}将此线程的\co{counter}加到\co{globalcount}上，
\clnref{glbrsv}从\co{globalreserve}中减去此线程的\co{countermax}。
\end{fcvref}

\QuickQuizSeries{%
\QuickQuizB{
	什么阻止线程在
	\cref{lst:count:Atomic Limit Counter Utility Functions 1}
	\clnrefr{ln:count:count_lim_atomic:utility1:flush:b}上的
	\co{flush_local_count()}清空其
	\co{counterandmax}变量后立即重新填充它？
}\QuickQuizAnswerB{
	在\co{flush_local_count()}的调用者释放
	\co{gblcnt_mutex}之前，其他线程无法重新填充其
	\co{counterandmax}。
	届时，\co{flush_local_count()}的调用者将已完成
	对计数值的使用，因此其他线程的重新填充不会有问题——
	假设\co{globalcount}的值足够大，允许重新填充。
}\QuickQuizEndB
%
\QuickQuizE{
	什么阻止\co{add_count()}或\co{sub_count()}快速路径的
	并发执行在
	\cref{lst:count:Atomic Limit Counter Utility Functions 1}
	\clnrefr{ln:count:count_lim_atomic:utility1:flush:atmxchg}上的
	\co{flush_local_count()}访问\co{counterandmax}变量时
	对其进行干扰？
}\QuickQuizAnswerE{
	没有什么能阻止。
	考虑以下三种情况：
	\begin{enumerate}
	\item	如果\co{flush_local_count()}的\co{atomic_xchg()}
		在任一快速路径的\co{split_counterandmax()}之前执行，
		则快速路径将看到零\co{counter}和\co{countermax}，
		从而转到慢速路径（除非\co{delta}为零）。
	\item	如果\co{flush_local_count()}的\co{atomic_xchg()}
		在任一快速路径的\co{split_counterandmax()}之后
		但在该快速路径的\co{atomic_cmpxchg()}之前执行，
		则\co{atomic_cmpxchg()}将失败，导致快速路径重启，
		这归结为上面的情况~1。
	\item	如果\co{flush_local_count()}的\co{atomic_xchg()}
		在任一快速路径的\co{atomic_cmpxchg()}之后执行，
		则快速路径（很可能）在\co{flush_local_count()}
		清零线程的\co{counterandmax}变量之前成功完成。
	\end{enumerate}
	无论如何，竞争都被正确解决。
}\QuickQuizEndE
}

\begin{fcvref}[ln:count:count_lim_atomic:utility2]
\cref{lst:count:Atomic Limit Counter Utility Functions 2}
\Clnrefrange{balance:b}{balance:e}
展示了\co{balance_count()}的代码，
它重新填充调用线程的本地\co{counterandmax}变量。
此函数与前面算法非常相似，
需要进行更改以处理合并的\co{counterandmax}变量。
代码的详细分析留作读者的练习，
\clnref{register:b}开始的\co{count_register_thread()}函数
和\clnref{unregister:b}开始的\co{count_unregister_thread()}函数
也同样留作练习。
\end{fcvref}

\QuickQuizSeries{%
\QuickQuizB{
	\cref{lst:count:Atomic Limit Counter Utility Functions 2}中
	\co{balance_count()}的
	\clnrefr{ln:count:count_lim_atomic:utility2:balance:atmcset}
	如何在并发\co{flush_local_count()}更新此变量的情况下
	正确工作？
}\QuickQuizAnswerB{
	如果它们并发执行的话，确实无法工作！
	但不用担心。
	\co{balance_count()}和\co{flush_local_count()}的调用者
	都持有\co{gblcnt_mutex}，因此在给定时刻只有一个
	可以执行。
}\QuickQuizEndB
%
\QuickQuizE{
	\co{balance_count()}中的\co{atomic_set()}原语
	真的需要是atomic的吗？
}\QuickQuizAnswerE{
	不需要。
	就相关代码而言，\co{atomic_set()}更新的是
	其自身线程的per-thread变量\co{counterandmax}，
	该变量的值只能由其自身线程的快速路径或
	所有更新方的慢速路径更新。
	由于\co{atomic_set()}在慢速路径中，它永远不会
	与其他atomic操作竞争。

	因此，即使此代码被移植到普通存储可以
	干扰对\co{atomic_t}的atomic读-修改-写操作的ISA上，
	非atomic的\co{atomic_set()}也应该能工作。

	话虽如此，这种仅对那些特殊ISA有效的
	慢速路径优化可能不值得费心。
}\QuickQuizEndE
}			% End of \QuickQuizSeries

下一节对此设计进行定性评估。
\ebFloatBarrier

\subsection{原子限值计数器讨论}

这是第一个真正允许counter一直增长到上限的实现，
但这也带来了快速路径上atomic操作的开销，
让快速路径明显变慢了。
虽然在某些场合这种变慢是允许的，
但仍然值得我们去探索让写入端性能更好的算法。
使用signal handler从其他线程窃取计数就是其中一种算法。
因为signal handler可以运行在收到信号的线程的上下文中，
所以就不需要atomic操作了，下一节将讨论这一算法。

\QuickQuiz{
	但signal handler可以在运行时被迁移到其他CPU上。
	这种可能性不需要atomic指令和内存屏障（memory barrier）来
	在线程和中断该线程的signal handler之间
	可靠通信吗？
}\QuickQuizAnswer{
	不需要。
	如果signal handler被迁移到另一个CPU，
	被中断的线程也随之迁移。
}\QuickQuizEnd

\subsection{信号窃取限值计数器设计}
\label{sec:count:Signal-Theft Limit Counter Design}

虽然per-thread状态只由对应线程更改，
但signal handler仍然有必要进行同步。
\cref{fig:count:Signal-Theft State Machine}所示的状态机
提供了这种同步。

\begin{figure}
\centering
\resizebox{2in}{!}{\includegraphics{count/sig-theft}}
\caption{Signal-Theft状态机}
\label{fig:count:Signal-Theft State Machine}
\end{figure}

状态机从IDLE状态开始，当\co{add_count()}或\co{sub_count()}
发现线程的本地计数和全局计数之和已经不足以容纳请求的大小时，
对应的慢速路径将每个线程的\co{theft}状态设置为REQ
（除非该线程没有计数值，这样它就直接转为READY）。
只有在慢速路径获得\co{gblcnt_mutex} lock后，才允许
从IDLE状态转换为其他状态，如绿色所示。\footnote{
	对于持有黑白版本本书的读者，
	IDLE和READY是绿色，REQ是红色，ACK是蓝色。}
然后慢速路径向每个线程发送一个信号，
对应的signal handler检查本线程的\co{theft}和
\co{counting}变量。
如果\co{theft}状态不为REQ，那么signal handler就不能
改变其状态，只能直接返回。
而\co{theft}状态为REQ时，如果设置了\co{counting}变量
（表明当前线程正处于快速路径），signal handler将\co{theft}状态
设置为ACK，而不是READY。

如果\co{theft}状态为ACK，
那么只有快速路径才有权改变\co{theft}状态，如蓝色所示。
当快速路径完成时，会将\co{theft}状态设置为READY。

一旦慢速路径发现某个线程的\co{theft}状态为READY，
这时慢速路径有权窃取此线程的计数。
然后慢速路径将线程的\co{theft}状态设置为IDLE。

\QuickQuizSeries{%
\QuickQuizB{
	在\cref{fig:count:Signal-Theft State Machine}中，
	为什么REQ \co{theft}状态用红色标记？
}\QuickQuizAnswerB{
	表示只有快速路径才允许改变\co{theft}状态，
	如果线程在此状态停留太久，
	运行慢速路径的线程将重新发送POSIX信号。
}\QuickQuizEndB
%
\QuickQuizE{
	在\cref{fig:count:Signal-Theft State Machine}中，
	拥有单独的REQ和ACK \co{theft}状态有什么意义？
	为什么不通过将它们合并为单个REQACK状态来简化状态机？
	那么signal handler或快速路径中先到达的
	就可以将状态设为READY。
}\QuickQuizAnswerE{
	合并REQ和ACK状态是一个非常糟糕的主意，
	原因包括：
	\begin{enumerate}
	\item	慢速路径使用REQ和ACK状态来确定
		是否应重新传输信号。
		如果状态被合并，慢速路径将别无选择
		只能发送冗余信号，这将产生
		不必要地减慢快速路径的不良效果。
	\item	将产生以下竞争：
		\begin{enumerate}
		\item	慢速路径将给定线程的状态设置为REQACK。
		\item	该线程刚刚完成其快速路径，
			并注意到REQACK状态。
		\item	线程接收到信号，signal handler也注意到
			REQACK状态，并且由于没有快速路径
			在进行中，将状态设为READY。
		\item	慢速路径注意到READY状态，窃取计数，
			将状态设为IDLE，并完成。
		\item	快速路径将状态设为READY，
			禁用了该线程的后续快速路径执行。
		\end{enumerate}
		这里的基本问题是合并的REQACK状态
		可以被signal handler和快速路径同时引用。
		四状态设置维护的清晰分离确保了有序的状态转换。
	\end{enumerate}
	话虽如此，你也许能够使三状态设置正确工作。
	如果你成功了，请仔细与四状态设置进行比较。
	三状态解决方案真的更好吗？为什么？
}\QuickQuizEndE
}

\subsection{信号窃取限值计数器实现}
\label{sec:count:Signal-Theft Limit Counter Implementation}

\begin{fcvref}[ln:count:count_lim_sig:data]
\Cref{lst:count:Signal-Theft Limit Counter Data}
（\path{count_lim_sig.c}）
展示了基于signal-theft的counter实现使用的数据结构。
\Clnrefrange{value:b}{value:e}定义了前一节描述的
per-thread theft状态机的状态和值。
\Clnrefrange{var:b}{var:e}与早期实现类似，
增加了\clnref{maxp,theftp}以允许远程访问
线程的\co{countermax}和\co{theft}变量。
\end{fcvref}

\begin{listing}
\ebresizeverb{.9}{\input{CodeSamples/count/count_lim_sig@data.fcv}}
\caption{Signal-Theft Limit Counter数据}
\label{lst:count:Signal-Theft Limit Counter Data}
\end{listing}

\begin{fcvref}[ln:count:count_lim_sig:migration:globalize]
\Cref{lst:count:Signal-Theft Limit Counter Value-Migration Functions}
展示了负责在per-thread变量和全局变量之间迁移计数的函数。
\Clnrefrange{b}{e}展示了\co{globalize_count()}，
与早期实现相同。
\end{fcvref}
\begin{fcvref}[ln:count:count_lim_sig:migration:flush_sig]
\Clnrefrange{b}{e}展示了\co{flush_local_count_sig()}，
它是theft过程中使用的signal handler。
\Clnref{check:REQ,return:n}检查\co{theft}状态是否为REQ，
如果不是则不做更改直接返回。
\Clnref{set:ACK}将\co{theft}状态设为ACK，
如果\clnref{check:fast}发现此线程的快速路径未在运行，
\clnref{set:READY}使用\co{smp_store_release()}
将\co{theft}状态设为READY，
进一步确保快速路径中对\co{counter}的任何更改
发生在\co{theft}更改为READY之前。
\end{fcvref}

\begin{listing}
\ebresizeverb{.83}{\input{CodeSamples/count/count_lim_sig@migration.fcv}}
\caption{Signal-Theft Limit Counter值迁移函数}
\label{lst:count:Signal-Theft Limit Counter Value-Migration Functions}
\end{listing}

\QuickQuiz{
	在\cref{lst:count:Signal-Theft Limit Counter Value-Migration Functions}中，
	\co{flush_local_count_sig()}不需要更强的内存屏障吗？
}\QuickQuizAnswer{
	不需要，那个\co{smp_store_release()}就够了，
	因为此代码仅与\co{flush_local_count()}通信，
	而且不需要store-to-load排序。
}\QuickQuizEnd

\begin{fcvref}[ln:count:count_lim_sig:migration:flush]
\Clnrefrange{b}{e}展示了\co{flush_local_count()}，
它从慢速路径调用以刷新所有线程的本地计数。
\clnrefrange{loop:b}{loop:e}的循环
为每个持有本地计数的线程推进\co{theft}状态，
并向该线程发送信号。
\Clnref{skip}跳过任何不存在的线程。
否则，\clnref{checkmax}检查当前线程是否持有本地计数，
如果没有，\clnref{READY}将线程的\co{theft}状态设为READY，
\clnref{next}跳到下一个线程。
否则，\clnref{REQ}将线程的\co{theft}状态设为REQ，
\clnref{signal}向线程发送信号。
\end{fcvref}

\QuickQuizSeries{%
\QuickQuizB{
	在\cref{lst:count:Signal-Theft Limit Counter Value-Migration Functions}中，
	为什么\clnrefr{ln:count:count_lim_sig:migration:flush:checkmax}
	可以安全地直接访问其他线程的\co{countermax}变量？
}\QuickQuizAnswerB{
	因为除非持有\co{gblcnt_mutex} lock，
	否则不允许其他线程更改其\co{countermax}变量的值。
	但调用者已获取此lock，因此其他线程不可能持有它，
	因此不允许更改其\co{countermax}变量。
	所以我们可以安全地访问它——但不能更改它。
}\QuickQuizEndB
%
\QuickQuizM{
	在\cref{lst:count:Signal-Theft Limit Counter Value-Migration Functions}中，
	为什么\clnrefr{ln:count:count_lim_sig:migration:flush:signal}
	不检查当前线程是否在向自己发送信号？
}\QuickQuizAnswerM{
	不需要额外检查。
	\co{flush_local_count()}的调用者已经调用了
	\co{globalize_count()}，因此
	\clnrefr{ln:count:count_lim_sig:migration:flush:checkmax}
	上的检查将成功，跳过后面的\co{pthread_kill()}。
}\QuickQuizEndM
%
\QuickQuizE{
	\cref{lst:count:Signal-Theft Limit Counter Data,%
	lst:count:Signal-Theft Limit Counter Value-Migration Functions}中
	所示的代码适用于\GCC\ 和POSIX。
	要使其同时符合ISO C标准需要什么？
}\QuickQuizAnswerE{
	\co{theft}变量必须是\co{sig_atomic_t}类型，
	以保证它可以在signal handler和
	被信号中断的代码之间安全共享。
}\QuickQuizEndE
}

\begin{fcvref}[ln:count:count_lim_sig:migration:flush]
\clnrefrange{loop2:b}{loop2:e}的循环
等待每个线程达到READY状态，然后窃取该线程的计数。
\Clnrefrange{skip:nonexist}{next2}跳过不存在的线程，
\clnrefrange{loop3:b}{loop3:e}的循环
等待当前线程的\co{theft}状态变为READY。
\Clnref{block}阻塞一毫秒以避免优先级反转问题，
如果\clnref{check:REQ}确定线程的信号尚未到达，
\clnref{signal2}重新发送信号。
当线程的\co{theft}状态变为READY时执行到达\clnref{thiev:b}，
因此\clnrefrange{thiev:b}{thiev:e}执行窃取。
\Clnref{IDLE}然后将线程的\co{theft}状态设回IDLE。
\end{fcvref}

\QuickQuiz{
	在\cref{lst:count:Signal-Theft Limit Counter Value-Migration Functions}中，
	为什么\clnrefr{ln:count:count_lim_sig:migration:flush:signal2}
	要重新发送信号？
}\QuickQuizAnswer{
	因为数十年来许多操作系统都有偶尔丢失信号的特性。
	这是特性还是bug尚有争议，但无关紧要。
	从用户的角度来看，明显的症状不是内核bug，
	而是用户应用程序挂起。

	\emph{你的}用户应用程序挂起！
}\QuickQuizEnd

\begin{fcvref}[ln:count:count_lim_sig:migration:balance]
\Clnrefrange{b}{e}展示了\co{balance_count()}，
与早期示例类似。
\end{fcvref}

\begin{listing}
\input{CodeSamples/count/count_lim_sig@add.fcv}
\caption{Signal-Theft Limit Counter加法函数}
\label{lst:count:Signal-Theft Limit Counter Add Function}
\end{listing}

\begin{listing}
\input{CodeSamples/count/count_lim_sig@sub.fcv}
\caption{Signal-Theft Limit Counter减法函数}
\label{lst:count:Signal-Theft Limit Counter Subtract Function}
\end{listing}

\begin{fcvref}[ln:count:count_lim_sig:add]
\Cref{lst:count:Signal-Theft Limit Counter Add Function}
展示了\co{add_count()}函数。
快速路径覆盖\clnrefrange{fast:b}{return:fs}，
慢速路径覆盖\clnrefrange{acquire}{return:ss}。
\Clnref{fast:b}将per-thread \co{counting}变量设为1，
这样任何后续中断此线程的signal handler将
把\co{theft}状态设为ACK而不是READY，
允许此快速路径正确完成。
\Clnref{barrier:1}防止编译器将快速路径主体的任何部分
重排到\co{counting}的设置之前。
\Clnref{check:b,check:e}检查per-thread数据是否能容纳
\co{add_count()}且没有正在进行的theft，
如果是，\clnref{add:f}执行快速路径加法，
\clnref{fasttaken}记录快速路径已被采用。

在两种情况下，\clnref{barrier:2}防止编译器将
快速路径主体重排到\clnref{clearcnt}之后，
后者允许任何后续signal handler进行theft。
\Clnref{barrier:3}再次禁用编译器重排，
然后\clnref{check:ACK}检查signal handler
是否延迟了\co{theft}状态到READY的更改，
如果是，\clnref{READY}使用\co{smp_store_release()}
将\co{theft}状态设为READY，
进一步确保任何看到READY状态的CPU
也看到\clnref{add:f}的效果。
如果\clnref{add:f}上的快速路径加法已执行，
则\clnref{return:fs}返回成功。
\end{fcvref}

\begin{listing}
\input{CodeSamples/count/count_lim_sig@read.fcv}
\caption{Signal-Theft Limit Counter读取函数}
\label{lst:count:Signal-Theft Limit Counter Read Function}
\end{listing}

\begin{fcvref}[ln:count:count_lim_sig:add]
否则，我们落入从\clnref{acquire}开始的慢速路径。
慢速路径的结构与早期示例类似，
其分析留作读者的练习。
\end{fcvref}
类似地，\cref{lst:count:Signal-Theft Limit Counter Subtract Function}上
\co{sub_count()}的结构与\co{add_count()}相同，
因此\co{sub_count()}的分析也留作读者的练习，
\cref{lst:count:Signal-Theft Limit Counter Read Function}中
\co{read_count()}的分析亦然。

\begin{listing}
\input{CodeSamples/count/count_lim_sig@initialization.fcv}
\caption{Signal-Theft Limit Counter初始化函数}
\label{lst:count:Signal-Theft Limit Counter Initialization Functions}
\end{listing}

\begin{fcvref}[ln:count:count_lim_sig:initialization:init]
\cref{lst:count:Signal-Theft Limit Counter Initialization Functions}的
\Clnrefrange{b}{e}
展示了\co{count_init()}，它将\co{flush_local_count_sig()}
设置为\co{SIGUSR1}的signal handler，
使\co{flush_local_count()}中的\apipx{pthread_kill()}调用
能够调用\co{flush_local_count_sig()}。
线程注册和注销的代码与早期示例类似，
其分析留作读者的练习。
\end{fcvref}

\subsection{信号窃取限值计数器讨论}

在我的六核x86笔记本电脑上，使用signal-theft的实现比
atomic操作的实现快八倍以上。
它总是这么好吗？

由于atomic指令的相对缓慢，signal-theft实现在Pentium-4处理器上
比atomic操作好得多，但后来，老式80386对称多处理器系统
在atomic操作实现的路径深度更短，atomic操作的性能也随之提升。
可是，更新端的性能提升是以读取端的高昂开销为代价的：
那些POSIX信号不是没有开销的。
如果考虑最终的性能，你需要在实际部署应用程序的
系统上测试这两种手段。

\QuickQuiz{
	POSIX信号不仅慢，向每个线程发送一个信号
	根本无法扩展。
	如果你有（比方说）10,000个线程，
	需要读取端快速，你会怎么做？
}\QuickQuizAnswer{
	一种方法是使用\cref{sec:count:Eventually Consistent Implementation}
	中展示的技术，在单个变量中汇总整体counter值的近似值。
	另一种方法是使用多个线程来执行读取，
	每个线程与更新线程的特定子集交互。
}\QuickQuizEnd

这就是为什么高质量的API如此重要的一个理由：
它们让实现代码可以随着持续变动的硬件性能特征一起改变。

\QuickQuiz{
	如果你想要一个exact limit counter仅对其下限精确，
	但允许上限不精确呢？
}\QuickQuizAnswer{
	一个简单的解决方案是将上限夸大所需的数量。
	这种夸大的极限情况是将上限设置为
	counter能够表示的最大值。
}\QuickQuizEnd

\subsection{精确限值计数器的应用}
\label{sec:count:Applying Exact Limit Counters}

虽然本节介绍的exact limit counter实现非常有用，
但如果counter的值总是在零附近变动，它们就帮不了什么忙了，
正如统计对I/O设备的未完成访问次数时可能出现的情况。
如果我们并不关心当前有多少计数，这种值总在零附近变动的
计数开销很大。
如\QuickQuizRef{\QcountQIOcnt}中提出的
可移除I/O设备访问计数问题所述，
除非有人想移除设备，否则访问次数完全不重要，
而移除设备这种情况本身又很少见。

一种简单的解决办法是给counter添加一个很大的``偏置值''
（例如十亿），确保counter的值远离零，让counter可以高效工作。
当有人想要移除设备时，再从counter值中减去偏置值。
计数最后几次访问将相当低效，
但对之前的许多访问却可以全速计数了。

\QuickQuiz{
	使用偏置counter时，你最好还做了什么？
}\QuickQuizAnswer{
	你最好已将上限设置得足够大，
	以容纳偏置、预期最大访问次数，
	以及足够的``余量''来允许counter即使在
	访问次数达到最大值时也能高效工作。
}\QuickQuizEnd

虽然偏置counter有效且有用，但它只是
\cpageref{chp:Counting}上提出的可移除I/O设备
访问计数问题的部分解决方案。
当尝试移除设备时，我们不仅需要知道当前精确的I/O访问计数，
还需要从现在开始阻止未来的访问请求。
一种解决办法是在更新counter时读获取读写lock，
在读取counter时写获取同一读写lock。
执行I/O的代码可能如下：

\begin{fcvlabel}[ln:count:inline:I/O]
\begin{VerbatimN}[commandchars=\\\[\]]
read_lock(&mylock);		\lnlbl[acq]
if (removing) {			\lnlbl[check]
	read_unlock(&mylock);	\lnlbl[rel1]
	cancel_io();		\lnlbl[cancel]
} else {
	add_count(1);		\lnlbl[inc]
	read_unlock(&mylock);	\lnlbl[rel2]
	do_io();		\lnlbl[do]
	sub_count(1);		\lnlbl[dec]
}
\end{VerbatimN}
\end{fcvlabel}

\begin{fcvref}[ln:count:inline:I/O]
\Clnref{acq}读获取lock，
\clnref{rel1}或~\lnref{rel2}释放它。
\Clnref{check}检查设备是否正在被移除，如果是，
\clnref{rel1}释放lock，
\clnref{cancel}取消I/O或采取设备即将移除时适当的操作。
否则，\clnref{inc}递增访问计数，
\clnref{rel2}释放lock，
\clnref{do}执行I/O，
\clnref{dec}递减访问计数。
\end{fcvref}

\QuickQuiz{
	这太荒谬了！
	我们在\emph{读}获取读写lock来\emph{更新}counter？
	你在搞什么？
}\QuickQuizAnswer{
	奇怪，也许，但确实如此！
	几乎足以让你觉得``读写lock''这个名字起得不好，不是吗？
}\QuickQuizEnd

移除设备的代码可能如下：

\begin{fcvlabel}[ln:count:inline:remove]
\begin{VerbatimN}[commandchars=\\\[\]]
write_lock(&mylock);		\lnlbl[acq]
removing = 1;			\lnlbl[note]
sub_count(mybias);
write_unlock(&mylock);		\lnlbl[rel]
while (read_count() != 0)	\lnlbl[loop:b]
	poll(NULL, 0, 1);	\lnlbl[loop:e]
remove_device();		\lnlbl[remove]
\end{VerbatimN}
\end{fcvlabel}

\begin{fcvref}[ln:count:inline:remove]
\Clnref{acq}写获取lock，\clnref{rel}释放它。
\Clnref{note}记录设备正在被移除，
\clnrefrange{loop:b}{loop:e}的循环等待所有I/O操作完成。
最后，\clnref{remove}执行为设备移除准备所需的额外处理。
\end{fcvref}

\QuickQuiz{
	在实际系统中还需要考虑哪些其他问题？
}\QuickQuizAnswer{
	非常多！

	以下是一些入门级的：

	\begin{enumerate}
	\item	可能有任意数量的设备，因此全局变量不合适，
		像\co{do_io()}这样的函数也缺少参数。
	\item	轮询循环在实际系统中可能有问题，
		浪费CPU时间和能量。
		在许多情况下，事件驱动设计要好得多，
		例如最后完成的I/O唤醒设备移除线程。
	\item	I/O可能失败，因此\co{do_io()}可能需要返回值。
	\item	如果设备故障，最后一个I/O可能永远不会完成。
		在这种情况下，可能需要某种超时来允许错误恢复。
	\item	\co{add_count()}和\co{sub_count()}都可能失败，
		但它们的返回值没有被检查。
	\item	读写lock扩展性不好。
		避免读写lock高读获取成本的一种方法在
		\cref{chp:Locking,chp:Deferred Processing}中介绍。
	\end{enumerate}
}\QuickQuizEnd

\section{并行计数讨论}
\label{sec:count:Parallel Counting Discussion}
%
\epigraph{具体中蕴含普遍性这一思想具有深远的重要性。}
	 {Douglas R. Hofstadter}

本章展示了传统计数原语会遇到的问题：可靠性、性能和可扩展性。
C语言的\co{++}运算符不能在多线程代码中保证可靠运行，
对单个变量的atomic操作性能不好，可扩展性也差。
因此，本章还展示了一系列在特殊情况下
性能和可扩展性俱佳的计数算法。

回顾这些计数算法的经验教训很有价值。
为此，\cref{sec:count:Parallel Counting Validation}
概述了必要的验证，
\cref{sec:count:Parallel Counting Performance}
总结了性能和可扩展性，
\cref{sec:count:Parallel Counting Specializations}
讨论了专门化的必要性，
最后，\cref{sec:count:Parallel Counting Lessons}
列举了所学的经验教训，并引导读者关注将在后续章节中
扩展这些教训的内容。

\begin{table*}
\rowcolors{4}{}{lightgray}
\renewcommand*{\arraystretch}{1.1}
\small
\centering
\newcommand{\NA}{\cellcolor{white}}
\ebresizewidth{
\begin{tabular}{lrcS[table-format=2.1]S[table-format=3.0]S[table-format=4.0]
		  S[table-format=6.0]S[table-format=6.0]}
	\toprule
	\multirow{2}{*}{\begin{picture}(60,15)(0,-3)\put(0,0){Algorithm}
			\put(14,-10){(\path{count_*.c})}\end{picture}} &
	    & \multirow{2}{*}{\begin{picture}(6,50)(0,-24)\rotatebox{90}{Exact?}\end{picture}} &
		\multicolumn{1}{c}{\multirow{2}{*}{\begin{picture}(30,15)(0,-3)
			\put(0,0){Updates}\put(15,-10){(ns)}\end{picture}}} &
			\multicolumn{4}{c}{Reads (ns)} \\
	\cmidrule{5-8}
	    & Section & & &
				   \multicolumn{1}{r}{1 CPU} &
				      \multicolumn{1}{r}{8 CPUs} &
					 \multicolumn{1}{r}{64 CPUs} &
					    \multicolumn{1}{r}{420 CPUs} \\
		\midrule
		\path{stat} & \ref{sec:count:Array-Based Implementation} & \NA &
		 6.3 & 294 & 303   & 315     &    612 \\
	\path{stat_eventual} & \ref{sec:count:Eventually Consistent Implementation} & \NA &
		 6.4 &   1 &   1   &   1     &      1 \\
	\path{end} & \ref{sec:count:Per-Thread-Variable-Based Implementation} & \NA &
		 2.9 & 301 & 6 309 & 147 594 & 239 683 \\
	\path{end_rcu} & \ref{sec:together:RCU and Per-Thread-Variable-Based Statistical Counters} & \NA &
		 2.9 & 454 &   481 &     508 &   2 317 \\
	\midrule
	\path{lim} & \ref{sec:count:Simple Limit Counter Implementation} &
		N &  3.2 & 435 & 6 678 & 156 175 & 239 422 \\
	\path{lim_app} & \ref{sec:count:Approximate Limit Counter Implementation} &
		N &  2.4 & 485 & 7 041 & 173 108 & 239 682 \\
	\path{lim_atomic} & \ref{sec:count:Atomic Limit Counter Implementation} &
		Y & 19.7 & 513 & 7 085 & 199 957 & 239 450 \\
	\path{lim_sig} & \ref{sec:count:Signal-Theft Limit Counter Implementation} &
		Y &  4.7 & 519 & 6 805 & 120 000 & 238 811 \\
	\bottomrule
\end{tabular}
}
\caption{x86上的Statistical/Limit Counter性能}
\label{tab:count:Statistical/Limit Counter Performance on x86}
\end{table*}

\subsection{并行计数验证}
\label{sec:count:Parallel Counting Validation}

本节中的许多算法相当简单，简单到令人倾向于
宣布它们构造正确或检查正确。
不幸的是，执行构造或检查的人太容易变得过于自信、
疲劳、困惑或只是粗心大意，所有这些都可能导致bug。
而且这些limit counter的早期实现确实包含了bug，
在某些情况下，在32位系统上维护64位计数的
固有复杂性助长了这些bug。
因此，即使对于本章中介绍的简单算法，验证也不是可选的。

Statistical counters通过充当counter的测试
（``\path{counttorture.h}''）进行验证，
即counter中的聚合总和随各个更新端线程添加的
金额总和而变化。

Limit counters也通过充当counter的测试
（``\path{limtorture.h}''）进行验证，
并额外检查其容纳指定限制的能力。

这两个测试套件都产生了
\cref{sec:count:Parallel Counting Performance}中使用的
性能数据。

虽然这种程度的验证对于像这样的教科书实现已经足够好，
但在将类似算法投入生产之前，应用额外的验证会是明智之举。
\Cref{chp:Validation}描述了测试的其他方法，
鉴于大多数这些计数算法的简洁性，
\cref{chp:Formal Verification}中描述的大多数技术也很有帮助。

\subsection{并行计数性能}
\label{sec:count:Parallel Counting Performance}

\cref{tab:count:Statistical/Limit Counter Performance on x86}的上半部分
展示了四种并行statistical counting算法的性能。
所有四种算法在更新方面都提供了近乎完美的线性可扩展性。
基于per-thread变量的实现（\path{count_end.c}）
在更新方面明显快于基于数组的实现
（\path{count_stat.c}），但在大量核心的读取上更慢，
并且在有许多并行读取者时遭受严重的\IX{lock contention}。
可以使用\cref{chp:Deferred Processing}中介绍的
延迟处理技术来解决这种竞争，
如\cref{tab:count:Statistical/Limit Counter Performance on x86}的
\path{count_end_rcu.c}行所示。
延迟处理在\path{count_stat_eventual.c}行上也表现出色，
这得益于最终一致性。

\QuickQuizSeries{%
\QuickQuizB{
	在\cref{tab:count:Statistical/Limit Counter Performance on x86}的
	\path{count_stat.c}行中，我们看到读取端随线程数量
	线性扩展。
	考虑到线程越多需要汇总的per-thread counter越多，
	这怎么可能？
}\QuickQuizAnswerB{
	读取端代码必须扫描整个固定大小的数组，
	不管线程数量多少，因此性能没有差异。
	相比之下，在后两种算法中，读取者在线程更多时
	必须做更多工作。
	此外，后两种算法引入了额外的间接层，
	因为它们需要从整数线程ID映射到
	相应的\co{_Thread_local}变量。
}\QuickQuizEndB
%
\QuickQuizE{
	即使在\cref{tab:count:Statistical/Limit Counter Performance on x86}的
	第四行，这些statistical counter实现的
	读取端性能也相当糟糕。
	那为什么还要使用它们？
}\QuickQuizAnswerE{
	``用合适的工具做合适的事。''

	从\cref{fig:count:Atomic Increment Scalability on x86}
	可以看出，单变量atomic递增
	不适合任何涉及大量并行更新的工作。
	相比之下，\cref{tab:count:Statistical/Limit Counter Performance on x86}
	上半部分所示的算法在处理更新密集的情况时
	表现出色。
	当然，如果你面对的是读取密集的情况，
	应该使用其他方法，例如具有单个atomic递增变量
	的最终一致性设计，该变量可以通过单次加载读出，
	类似于\cref{sec:count:Eventually Consistent Implementation}
	中使用的方法。
}\QuickQuizEndE
}

\cref{tab:count:Statistical/Limit Counter Performance on x86}的下半部分
展示了并行limit counting算法的性能。
对上限计数的精确性要求带来了不小的更新端性能损失，
不过在这个x86系统上，用信号替代atomic操作
可以减轻这种损失。
所有这些实现在面对多个读取者并发读取时都存在
读取端lock竞争问题。

\QuickQuizSeries{%
\QuickQuizB{
	根据\cref{tab:count:Statistical/Limit Counter Performance on x86}
	下半部分所示的性能数据，
	我们应该总是优先选择信号而不是atomic操作，对吧？
}\QuickQuizAnswerB{
	这取决于工作负载。
	注意在64核系统上，你需要超过一百次非atomic操作
	（每次大约有40纳秒的性能增益）才能弥补
	一次信号的开销（接近5\emph{微秒}的性能损失）。
	虽然不乏具有更高读取强度的工作负载，
	你需要考虑你的特定工作负载。

	此外，虽然内存屏障在历史上与普通指令相比一直很昂贵，
	你应该在将要运行的特定硬件上检查这一点。
	计算机硬件的特性确实会随时间变化，
	算法也必须相应地变化。
}\QuickQuizEndB
%
\QuickQuizE{
	是否可以应用高级技术来解决
	\cref{tab:count:Statistical/Limit Counter Performance on x86}
	下半部分中读取者看到的lock竞争？
}\QuickQuizAnswerE{
	一种方法是牺牲一些更新端性能，
	就像可扩展非零指示器
	（SNZI）~\cite{FaithEllen:2007:SNZI}所做的那样。
	还有其他多种方法可以做到这一点，
	留作读者的练习。
	任何应用层次结构的方法——用对应于层次结构较低级别的
	本地lock获取替代频繁的全局lock获取——都应该效果很好。
}\QuickQuizEndE
}

简而言之，本章使用的算法只在它们特定的领域运转良好。
但我们的并行计数算法是不是只能局限于这些特殊场合呢？
如果有一种通用算法能在所有场合都性能出色是不是更好？
下一节将探讨这些问题。

\subsection{并行计数的专门化}
\label{sec:count:Parallel Counting Specializations}

上述算法只在各自的领域性能出色，
这可以说是并行计算的一个主要问题。
毕竟，C语言的\co{++}运算符在所有单线程程序中性能都不错，
不仅仅是个别领域，而是所有情况，对吧？

上述推论有一点道理，但本质上是误导的。
我们提到的问题不仅是并行性，更是可扩展性。
要弄明白这个，先来看看C语言的\co{++}运算符。
事实上\co{++}运算符并不是``总能''工作的，
只对一定范围的数字而言有效。
假如你要处理1,000位的十进制数，
C语言的\co{++}运算符对你来说就没用了。

\QuickQuiz{
	\co{++}运算符对1,000位数完全可以工作！
	你没听说过运算符重载吗？
}\QuickQuizAnswer{
	在C++语言中，你也许确实可以对1,000位数使用\co{++}，
	假设你能访问一个实现此类数字的类。
	但截至2021年，C语言不允许运算符重载。
}\QuickQuizEnd

我们提到的问题也不专属于算术。
假设你需要存储和查询数据，
是不是还会用ASCII文件、XML、关系型数据库、链表、
密集数组、B树、基数树或者其他什么数据结构和环境来存取数据？
这取决于你需要做什么、需要它做得有多快、以及数据集有多大
——即使在串行系统上也是如此。

同样，如果你需要计数，你的方案取决于统计的数有多大、
有多少个CPU并发操纵计数、如何使用计数，
以及所需要的性能和可扩展性的程度。

这个问题也不仅限于软件。
只需要简单地铺一块木板，就是一座让人跨过小溪的桥。
但你不能用一块木板横跨哥伦比亚河数公里宽的出海口，
也不能用于承载混凝土运输车的桥梁。
简而言之，桥的设计必须随着跨度和负载而改变，
并且，软件也要能适应硬件或者工作负载的变化，
能够自动进行最好。
事实上已经有一些关于这方面自动化的研究
~\cite{Appavoo03a,Soules03a}，
Linux内核也做了一些启动时的动态设定，
其中包括有限的二进制重写。
随着主流系统上CPU个数的逐步增加，
软件的动态适应变得越来越重要。

正如\cref{chp:Hardware and its Habits}所讨论的，
物理定律对并行软件的限制并不逊于对桥梁等
机械制品的限制。
这些限制迫使软件针对特定问题特定处理，
不过软件还是可以动态地选择特定算法，
以适应不同情况下的硬件和工作负载。

当然，即使是通用计数也是相当专门化的。
我们需要用计算机做大量其他事情。
下一节将我们从counter中学到的东西
与本书后续章节讨论的主题联系起来。

\subsection{并行计数的经验教训}
\label{sec:count:Parallel Counting Lessons}

本章开头承诺读者，对计数的研究是并行编程的绝佳切入点。
本节把本章学到的经验教训对应到后续章节的内容。

本章的例子表明，\emph{分区}是提升可扩展性和性能的重要工具。
计数有时可以完全被分区，
如\cref{sec:count:Statistical Counters}中的statistical counters，
或者部分被分区，
如\cref{sec:count:Approximate Limit Counters,%
sec:count:Exact Limit Counters}中的limit counters。
分区将在\cref{chp:Partitioning and Synchronization Design}中
更深入地讨论，
特别是部分并行化将在\cref{sec:SMPdesign:Parallel Fastpath}中讨论，
在那里它被称为\emph{parallel fastpath}。

\QuickQuiz{
	但如果我们要将所有东西都分区，
	为什么还要使用共享内存多线程？
	为什么不完全分区问题并作为多个进程运行，
	每个进程在自己的地址空间中？
}\QuickQuizAnswer{
	确实，拥有独立地址空间的多进程可以是
	利用并行性的绝佳方式，
	fork-join方法论和Erlang语言的支持者会非常
	乐意告诉你这一点。
	然而，共享内存并行性也有一些优势：
	\begin{enumerate}
	\item	只有应用程序中性能最关键的部分
		需要被分区，这些部分通常是
		应用程序的一小部分。
	\item	虽然cache miss与单个寄存器到寄存器指令
		相比相当慢，但它们通常比进程间通信原语
		快得多，而进程间通信原语反过来又比
		TCP/IP网络之类的东西快得多。
	\item	共享内存多处理器容易获得且价格低廉，
		因此，与1990年代形成鲜明对比的是，
		使用共享内存并行性几乎没有成本惩罚。
	\end{enumerate}
	一如既往，用合适的工具做合适的事！
}\QuickQuizEnd

部分分区的计数算法用lock来保护全局数据，
而lock正是\cref{chp:Locking}的主题。
与之相反，分区后的数据一般都完全处于各个线程的控制，
所以无须使用同步。
这种\emph{data ownership}将在
\cref{sec:SMPdesign:Data Ownership}中介绍，
并在\cref{chp:Data Ownership}中得到更详细的讨论。

由于整数加法和减法与典型的同步操作相比极其廉价，
要实现合理的可扩展性，需要谨慎使用同步操作。
实现这一目标的一种方法是批量处理加法和减法操作，
使大量这些廉价操作由单个同步操作处理。
\cref{tab:count:Statistical/Limit Counter Performance on x86}中
列出的每种计数算法都使用了某种形式的批量优化。

最后，\cref{sec:count:Eventually Consistent Implementation}中
提到的最终一致性statistical counter展示出，
延迟处理（在这里是指更新全局counter）
可以显著提升性能和可扩展性。
这种方法允许常见情况下的代码使用
比其他方式便宜得多的同步操作。
\Cref{chp:Deferred Processing}将研究很多可以提升性能、
可扩展性甚至实时响应速度的延迟处理方法。

总结要点：

\begin{enumerate}
\item	分区能够提升性能和可扩展性。
\item	部分分区——也就是只分区主要情况的代码路径——
	性能也很出色。
\item	部分分区可以应用在代码上（如
	\cref{sec:count:Statistical Counters}的statistical
	counters只分区了写操作，没有分区读操作），
	但是也可以应用在时间上
	（如\cref{sec:count:Approximate Limit Counters}和
	\cref{sec:count:Exact Limit Counters}的limit counters
	在离上限较远时运行很快，离上限较近时运行变慢）。
\item	跨时间分区通常在本地批量更新，
	以减少昂贵的全局操作次数，
	从而减少同步开销，进而改善性能和可扩展性。
	\cref{tab:count:Statistical/Limit Counter Performance on x86}
	中展示的所有算法都大量使用了批量处理。
\item	读取端的代码路径应保持只读：
	对共享内存的无谓同步写入会严重降低性能和可扩展性，
	如\cref{tab:count:Statistical/Limit Counter Performance on x86}的
	\path{count_end.c}行所示。
\item	经过审慎思考后的延迟处理能够提升性能和可扩展性，
	如\cref{sec:count:Eventually Consistent Implementation}所示。
\item	并行性能和可扩展性通常是跷跷板的两端：
	到达某种程度后，对代码的优化反而会降低另一方的表现。
	\cref{tab:count:Statistical/Limit Counter Performance on x86}
	的\path{count_stat.c}和\path{count_end_rcu.c}行
	说明了这一点。
\item	对性能和可扩展性的不同需求，以及其他很多因素，
	会影响算法和数据结构的设计。
	\Cref{fig:count:Atomic Increment Scalability on x86}
	展现了这一点：
	atomic递增对于双CPU系统来说完全可以接受，
	但对于八CPU系统来说就完全不合适了。
	% @@@ Rework the above sentence.
\end{enumerate}

\begin{figure}
\centering
\resizebox{3in}{!}{\includegraphics{count/FourTaskOrderOpt}}
\caption{优化与四大并行编程任务}
\label{fig:count:Optimization and the Four Parallel-Programming Tasks}
\end{figure}

进一步总结，我们有提高性能和可扩展性的``三大''方法，即
(1)~跨CPU或线程\emph{分区}，
(2)~\emph{批量处理}使每次昂贵的同步操作能完成更多工作，
(3)~在可行时\emph{弱化}同步操作。
作为粗略的经验法则，你应该按此顺序应用这些方法，
正如之前在讨论
\cpageref{fig:intro:Ordering of Parallel-Programming Tasks}上的
\cref{fig:intro:Ordering of Parallel-Programming Tasks}时所提到的。
分区优化适用于``资源分区和复制''气泡，
批量优化适用于``工作分区''气泡，
弱化优化适用于``并行访问控制''气泡，
如\cref{fig:count:Optimization and the Four Parallel-Programming Tasks}
所示。
当然，如果你使用专用硬件，如数字信号处理器（DSP）、
现场可编程门阵列（FPGA）或通用图形处理单元（GPGPU），
你可能需要在整个设计过程中密切关注
``与硬件交互''气泡。
例如，GPGPU硬件线程和内存连接的结构
可能极大地回报非常仔细的分区和批量处理设计决策。

简而言之，正如本章一开始介绍的，计数问题所蕴含的概念虽然简单，
却让我们得以探索许多并发方面的基础问题，
而不用分神于复杂的数据结构或精巧的同步手段。
后续的章节会对这些基础问题进行更深入的研究。

\QuickQuizAnswersChp{qqzcount}
